\chapter{Monoidal E-Graphs}
\label{ch:monoidal}

Building up on the insights from the previous chapter, we will present a generalisation of e-graphs tailored for more general symmetric monoidal theories.
We saw in~\cref{chapter:introduction} that treating symmetric monoidal theories algebraically and then applying e-graphs for equality saturation, i.e., for computing the members of the quotient algebra, is infeasible.
The proposal is to use string diagrams (or rather their concrete combinatorial representation as hypergraphs) and equip them with the ability to represent equivalent morphisms as a single object, akin to e-graphs.
We characterise such string diagrams as morphisms in a free semilattice-enriched symmetric monoidal category.
We then represent these string diagrams concretely as cospans of hierarchical hypergraphs and accordingly adapt the notion of DPOI rewriting for them.
Traditional e-graph can then be matched with string diagrams for a free \emph{cartesian} semilattice-enriched categories---a semilattice-enriched version of a Lawvere theory.
We saw in~\cref{chapter:algebras} that e-graphs can be seen as a device to compute the quotient for the hom-sets of a Lawvere theory.
Here, we characterise e-graphs as morphisms.
E-graph operations, such as \texttt{add} and \texttt{merge}, are then interpreted as string diagram rewriting.

\begin{remark}
The results in this chapter form a reworked version of the results presented in~\cite{tiurin2025equivalencehypergraphsdporewriting}.
\end{remark}

\section{$\catname{SLat}$-PROPs}

Let $\widetilde{F}_{\catname{SLat}}$ be the $2$-functor
$\catname{Cat} \to \catname{SLat}\text{-}\catname{Cat}$ induced by the free semilattice
functor $F_{\catname{SLat}} : \catname{Set} \to \catname{SLat}$
of~\cref{chapter:enriched}.

\begin{definition}\label{def:slat-prop}
An \emph{$\catname{SLat}$-PROP} is an $\catname{SLat}$-enriched strict symmetric monoidal
category with objects the natural numbers and $\otimes$ given by addition.
We write $\catname{PROP}^{\join}$ for the category of $\catname{SLat}$-PROPs and
identity-on-objects strict symmetric monoidal $\catname{SLat}$-functors.
\end{definition}

\begin{proposition}\label{prop:prop-adjunction}
The adjunction of~\cref{cor:2-adjunction-monoidal} restricts to an adjunction
\[
\widetilde{F}_{\catname{SLat}}, \widetilde{U}_{\catname{SLat}} :
\catname{PROP} \rightleftarrows \catname{PROP}^{\join}
\qquad \widetilde{F}_{\catname{SLat}} \dashv \widetilde{U}_{\catname{SLat}}~.
\]
\end{proposition}
\begin{proof}
Both $2$-functors are the identity on objects of the underlying categories, so they preserve
the object set $\mathbb{N}$, and they carry $\otimes$ to $\otimes$
by~\cref{ex:free-enrichment-monoidal}.
Strictness is preserved because the associators, unitors and symmetries of a PROP are
identity $2$-cells and a $2$-functor preserves these; for the symmetry this uses that
$F_{\catname{SLat}}$ is symmetric monoidal (\cref{cor:slat-free-monoidal}).
On $1$-cells, neither $2$-functor changes the action on objects, so identity-on-objects is
preserved in both directions, and the structure maps of a transpose $\overline{T}$ are
determined by their values on the injections $\iota_{u}$, hence are identities precisely when
those of $T$ are.
Conversely, if $\overline{T}$ is identity-on-objects and strict then so is
$T = \widetilde{U}_{\catname{SLat}}(\overline{T}) \circ \psi_{\mathcal{P}}$, since
$\psi_{\mathcal{P}}$ is identity-on-objects and strict monoidal.
As the morphisms of $\catname{PROP}$ and $\catname{PROP}^{\join}$ admit no non-identity
$2$-cells, the $2$-adjunction reduces to an ordinary adjunction.
\end{proof}

\begin{definition}
Let $\mathbf{S}_{\Sigma}^{\join}$ and $\mathbf{S}_{\Sigma, \mathcal{E}}^{\join}$ be categories given by the action of $\widetilde{F}$ on objects $\mathbf{S}_{\Sigma}$ and $\mathbf{S}_{\Sigma, \mathcal{E}}$ of $\catname{PROP}$.
\end{definition}

Equivalently, $\mathbf{S}_{\Sigma}^{\join}$ can be defined as an $\catname{SLat}$-category with objects natural numbers and hom-objects free semilattices given by sets freely constructed from the morphisms generated by the string diagrams in~\cref{fig:egraph-strings} modulo the equation of a symmetric monoidal category and the equations in~\cref{fig:string-equations}.
We extend the traditional syntax of string diagrams as seen in~\cref{chapter:string_diagrams} with a hierarchical dashed-box construct, inspired by functorial boxes of Mellies~\cite{mellies_functorial_2006}, to intuitively represent morphisms which are deemed equivalent---vertical dashed lines delimit the equivalent components of the box akin to an e-class in an e-graph.
Technically, the box just represents the semilattice structure on the hom-objects---delimiting lines correspond to $\join$-ing the enclosed morphisms. 
Even though the box is not functorial in the sense that it does not correspond to a functor, it enjoys similar \enquote{topological} properties as functorial boxes: one can slide other morphisms in and out of the box, provided they are duplicated and unduplicated accordingly.
\begin{figure}
\[
	\scalebox{0.65}{
	\tikzfig{./figures/monoidal_egraphs/egraph-strings}
	}
\]
\caption{String diagrams for semilattice-enriched symmetric monoidal categories.}
\label{fig:egraph-strings}
\end{figure}
\begin{figure}
	\[  
		\scalebox{0.6}{
		\tikzfig{./figures/monoidal_egraphs/egraph-strings-equations}
		}
	\]
	\caption{Equations for a  semilattice-enriched symmetric monoidal category.}
	\label{fig:string-equations}
\end{figure}
Similar to the cases of $\mathbf{H}_{\Sigma}$ and $\mathbf{S}_{\Sigma}$, we can define rewriting on $\mathbf{S}_{\Sigma}^{\join}$
\begin{definition}
A rewriting system $\mathcal{R}$ on an enriched PROP $\mathbb{A}$ consists of a set of rewrite rules, i.e, pairs $\langle l, r \rangle$ of morphisms $l, r : i \to j$ in $\mathbb{A}$ with the same domain and codomain.
Given $a,b : m \to n \in \mathbb{A}$, we say that $a$ rewrites to $b$ via $\mathcal{R}$, written $a \Rightarrow_{\mathcal{R}} b$, if either they are decomposable as follows
 \[
\tikzfig{./figures/string/decomposition}~,
 \] 
 or as
 \[
 \tikzfig{./figures/monoidal_egraphs/enriched_decomposition}~,
 \]
 such that $a' \Rightarrow_{\mathcal{R}} b'$. 
\end{definition}

This is already sufficient to give string diagrammatic interpretation to e-graphs.
\begin{example}\label{example:egg_egraph}
Consider $\mathbf{S}_{\Sigma, \mathcal{E}}$ additionally equipped with a natural commutative comonoid structure given by generators $\{\adjustbox{scale=0.5}{\tikzfig{./figures/monoidal_egraphs/copy_generator}} : 1 \to 2, \adjustbox{scale=0.5}{\tikzfig{./figures/monoidal_egraphs/delete_generator}} : 1 \to 0\}$ and equations of a commutative comonoid
\[
\adjustbox{width=\textwidth}{\tikzfig{./figures/monoidal_egraphs/cartesian_equations}}~.
\]
These are comonoid part of a commutative Frobenius algebra with the additional requirement of naturality for both the comultiplication and the counit.
Let us denote the resulting PROP as $\mathbf{Q}_{\Sigma, \mathcal{E}}$.
By Fox's theorem~\ref{theorem:fox}, this turns $\mathbf{Q}_{\Sigma, \mathcal{E}}$ into a cartesian category and makes it equivalent to the Lawvere theory $\mathbf{L}_{\Sigma, \mathcal{E}}$.
Applying $\widetilde{F}$ additionally equips the hom-sets with a semilattice structure yielding $\mathbf{Q}_{\Sigma, \mathcal{E}}^{\join}$.

Let 
\[
\Sigma_{1} = \{a : 0 \to 1, 2 : 0 \to 1, 1 : 0 \to 1, + : 2 \to 1, * : 2 \to 1, / : 2 \to 1, <\!\!< : 2 \to 1\}
\]
and
\begin{align*}
\mathcal{E} = \{x * 2 &= x <\!\!< 1\\
(x * y) / z &= x * (y / z)\\
(x / x) &= 1\}\\
\end{align*}

This is an algebraification of the~\cref{example:egraph_example_introduction} that we saw in the introduction.
Then consider the initial term $(a * 2) / 2$.
The e-graph for this term is given in~\cref{fig:egg_egraph} (a).
Each e-class of e-graph (a) contains a single e-node as none of the equations has been applied yet and noting that e-class containing the e-node $2$ is shared to avoid duplication.
E-graph (b) is obtained from (a) by applying the equation $x * 2 = x <\!\!< 1$ (from left to right).
As we saw in~\cref{chapter:algebras}, applying this equations amounts to finding a match within the e-graph (a), in this case $x$ is matched with the middle e-class containing $*$, instantiating the right-hand side of the rule with the obtained substitution, adding the right-hand side to the e-graph and merging the e-classes corresponding to the left- and righ-hand sides of the rule.
Graphically this can be depicted as. 
\[
\adjustbox{scale=0.8}{
\tikzfig{./figures/monoidal_egraphs/e-graph-example-a-no-label}
}
\xRightarrow{\texttt{add}(a <\!\!< 1)}
\adjustbox{scale=0.8}{
\tikzfig{./figures/monoidal_egraphs/e-graph-example-b-add}
}
\xRightarrow{\texttt{merge}(\scalebox{0.3}{\begin{tikzpicture}\begin{pgfonlayer}{nodelayer}\node [style=empty diag yellow] (0) at (0, 0) {<\!\!<};\end{pgfonlayer}\end{tikzpicture}},\scalebox{0.3}{\begin{tikzpicture}\begin{pgfonlayer}{nodelayer}\node [style=empty diag black] (0) at (0, 0) {*};\end{pgfonlayer}\end{tikzpicture}})}
\adjustbox{scale=0.8}{
\tikzfig{./figures/monoidal_egraphs/e-graph-example-b-merged}
}
\]
Then, applying the equation $(x * y) / z = x * (y / z)$ (again, from left to right) results in the following
\[
\scalebox{0.7}{
\tikzfig{./figures/monoidal_egraphs/e-graph-example-b-merged}
}
\xRightarrow{\texttt{add}(a * (2/2))}
\scalebox{0.7}{
\tikzfig{./figures/monoidal_egraphs/e-graph-example-c-add}
}
\xRightarrow{\texttt{merge}(\scalebox{0.3}{\begin{tikzpicture}\begin{pgfonlayer}{nodelayer}\node [style=empty diag yellow] (0) at (0, 0) {/};\end{pgfonlayer}\end{tikzpicture}},\scalebox{0.3}{\begin{tikzpicture}\begin{pgfonlayer}{nodelayer}\node [style=empty diag black] (0) at (0, 0) {*};\end{pgfonlayer}\end{tikzpicture}})}
\scalebox{0.7}{
\tikzfig{./figures/monoidal_egraphs/e-graph-example-c-merged}
}
\]
The resulting e-graph, after we applied the merge, may not always be valid, as it may violate the invariants and therefore additional steps may be necessary, like merging upwards to ensure congruence closure, or recanonicalising the e-nodes.
The latter is done implicitly in the picture: as can be noticed, every e-node has a single \emph{parent} e-class.
The other steps of e-graph transformation are similar and may be found in \textcolor{red}{TODO: appendix}.

We will now show that the same structure can be expressed using string diagrams for morphisms in $\mathbf{Q}_{\Sigma, \mathcal{E}}^{\join}$. 
The equivalent string diagram for each of the e-graphs is shown in~\cref{fig:egg_egraph_string}.
In the case of e-graph (a), the difference is almost cosmetical: instead of connecting e-nodes to e-classes, we connect the dashed boxes with each other and share the box with node $2$ using the comultiplication $\adjustbox{scale=0.3}{\tikzfig{./figures/monoidal_egraphs/copy_generator}}$.
We chose to wrap single nodes in boxes to make the correspondence more clear, but strictly speaking, this is not necessary and we will disallow such wrapped singletons and will write them without the box to avoid ambiguity between string diagrams and terms.
The string diagram for e-graph (b) is a little bit more interesting: the middle dashed box has now 2 components delimited by the vertical dashed line.
As boxes connected to each other, rather than nodes being connected to other nodes inside the boxes (that would break string diagram boundaries), we parameterise the whole box by a shared context and then components choose what to take from that context by explicitly using delete nodes $\adjustbox{scale=0.3}{\tikzfig{./figures/monoidal_egraphs/delete_generator}}$.
One can collapse all boxes using the equations of a semilattice-enriched category to get the set of terms represented by the string diagram.
The transformation from (a) to (b) is then done by the following rewriting of string diagrams:
\[
\adjustbox{width=\textwidth}{
    \tikzfig{./figures/monoidal_egraphs/egraph-translation-step-by-step-a-b}
}
\]
First, we duplicate the node $2$ using the naturality of the comultiplication, then we duplicate $a * 2$ in a box using idempotency of $\join$ and apply the equation to one of the components.
Then we unify the context of the two components using weakening and pull this context out of the box and re-share the node $2$, again, using naturality.
The resulting string diagram has the same degree of sharing as the corresponding e-graph (b).
The rest steps can be found in \textcolor{red}{TODO:appendix}.
\end{example}

\begin{figure}
\adjustbox{width=\textwidth}{
\tikzfig{./figures/monoidal_egraphs/egg_egraph}
}
\caption{E-graphs for the algebraic theory in~\cref{example:egg_egraph}.}
\label{fig:egg_egraph}
\end{figure}

\begin{figure}
\adjustbox{width=\textwidth}{
\tikzfig{./figures/monoidal_egraphs/egraph-translation}
}
\caption{String diagrams corresponding to e-graphs in~\cref{fig:egg_egraph}.}
\label{fig:egg_egraph_string}
\end{figure}

To make the string diagrams not a mere notation for terms, we will similarly turn to a category of hypergraphs.
This time, hypergraphs will need to contain sub-hypergraphs corresponding to the elements of the join $\join$.
For this, we will take an approach of hierarchical hypergraphs that appeared in the literature in various forms~\cite{Gaducci:hierarchical-graphs,plump:hierarchical-graphs,palacz:hierarchical-transform,fscd}.
We will call our version of hierarchical hypergraphs rooted \emph{e-hypergraphs} with \emph{e} standing for equivalence akin to e-graphs.

\section{Rooted E-Hypergraphs}
%% ---------------------------------------------------------------------------
%% Rooted e-hypergraphs
%% ---------------------------------------------------------------------------
\begin{remark}
Below we will elide the injections into the coproduct $\coprod$.
That is, given two objects $A$ and $B$ and a function $f : A \coprod B \to C$ or a relation $R \subseteq A \coprod B$ and elements $x : A$ and $y : B$, we will write $f(x,y)$ and $x R y$ instead of $f(\iota_A(x), \iota_B(y))$ and $\iota_A(x) R \iota_B(y)$.
We will also use the set union $\cup$ assuming that the sets of interest are disjoint.
\end{remark}
\begin{definition}[Rooted e-hypergraph]\label{def:rooted_ehypergraph}
A \emph{rooted e-hypergraph} $\mathcal{G}$ over a monoidal signature $\Sigma$ is a tuple
\[
(V, E, s, t, l, \top, <^{\mu}),
\]
where $(V, E, s, t)$ is a finite unlabelled hypergraph~\cref{def:hypergraph}, $l : E \to \Sigma_1 \cup \{\boxsort, \cmpsort\}$
is a labelling function, $\top \in E$, and $<^{\mu}$ is the immediate predecessor relation such that $<$ is a strict partial order on $V \coprod E$ given by a transitive closure of $<^{\mu}$ called the \emph{hierarchical relation}.
We write $E_{\Sigma}$, $E_{\boxsortsubscript}$ and $E_{\cmpsortsubscript}$ for the preimages of
$\Sigma_1$, $\boxsort$ and $\cmpsort$ under $l$, and call their elements
\emph{labelled}, \emph{box} and \emph{component} edges respectively.
We write $[x) \defeq \{x' \mid x' < x\}$ for the set of parents of $x$ and
$e <^{\mu} x$ if $e$ is the immediate predecessor of $x$.
We will also use terms predecessor and parent as well as successor and child interchangeably.
We require the following conditions.
\begin{enumerate}
  \item \emph{(Typing)} For $e \in E_{\Sigma}$, $|s(e)| = m$ and $|t(e)| = n$ where
        $l(e) : m \to n$; for $e \in E_{\cmpsortsubscript}$, $s(e) = t(e) = \varepsilon$.
  \item \emph{(Tree structure)} $\top \in E_{\cmpsortsubscript}$, $[\top) = \varnothing$, every
        $x \neq \top$ has exactly one immediate parent; consequently $(V \coprod E, <^{\mu})$ is a tree with
        root $\top$.
  \item \emph{(Alternation)} The immediate parent of a vertex, a labelled edge ($e \in E_{\Sigma}$) or a box
        edge ($e \in E_{\boxsortsubscript}$) is a component edge ($e \in E_{\cmpsortsubscript}$); the immediate parent of a component edge other
        than $\top$ is a box edge.
  \item \emph{(Incidence)} If $v \in s(e)$ or $v \in t(e)$ for
        $e \in E_{\Sigma} \cup E_{\boxsortsubscript}$, then $v$ and $e$ have the same
        immediate parent.
  \item \emph{(Non-degeneracy)} Every box edge has at least two children.
\end{enumerate}
We write $\mathfrak{T}$ for the rooted e-hypergraph
$(\varnothing, \{\top\}, \top \mapsto \varepsilon, \top \mapsto \varepsilon, l(\top) = \cmpsort, \top, \varnothing)$.
\end{definition}

\begin{figure}
  \[
  \adjustbox{max width=\textwidth}{
  \tikzfig{./figures/monoidal_egraphs/rooted-e-hypergraph-example}
  }
  \]
  \caption{Rooted e-hypergraph example.}
  \label{fig:rooted_ehypergraph_example}
\end{figure}

\begin{example}
Consider the rooted e-hypergraph shown in~\cref{fig:rooted_ehypergraph_example} for a signature $\Sigma = (\{*\}, \{ f : 1 \to 1, g : 1 \to 2\})$.
We depict edges from $E_{\Sigma}$ as round boxes labeled by the signature symbol, edges from $E_{\cmpsortsubscript}$ as boxes filled with gray, and edges from $E_{\boxsortsubscript}$ as boxes with dashed frame.
Sources are connected to the bottom side of the corresponding edge and targets are connected to the top side.
In this particular example we have
\[
\begin{array}{cccc}
  \top <^{\mu} v_1, v_2, w_1 & e_1 <^{\mu} e_2, e_3 & e_2 <^{\mu} u_1, u_2, z_1 & e_3 <^{\mu} u_3, z_2, z_3\\
  \top <^{\mu} e_1 & & e_2 <^{\mu} f &  e_3 <^{\mu} g~.
\end{array}
\]
We will occasionally abuse the labels of edges to use them as references to the said edges.
\end{example}

\begin{remark}[Comparison with~\cite{tiurin2025equivalencehypergraphsdporewriting}]\label{remark:compare_ehypergraph}
\Cref{def:rooted_ehypergraph} refines the e-hypergraphs
of~\cite[Def.~IV.2]{tiurin2025equivalencehypergraphsdporewriting} on multiple accounts.
There, the hierarchy was captured by a strict partial order in which an
element has \emph{at most} one immediate parent, so that parent-less
(\enquote{top-level}) elements coexist at an implicit outermost layer; here
the hierarchy is a genuine tree, with a single root $\top$ adjoined and every
other element having exactly one parent, so that \enquote{top-level} becomes
the definable notion \enquote{child of $\top$}, and $T$ replaces the empty
e-hypergraph as the smallest object.
Second, the components of an e-box were separated by a primitive
\emph{consistency relation} $\smile_{p}$, parameterised by a hierarchical hyperedge $p$---an equivalence relation on each sibling
set, required to be closed under connectivity and to have at least two
classes; here components are reified as $\cmpsort$-labelled edges,
consistency is \emph{derived}---$x \smile_{p} y$ precisely when there exist hyperedges $e_1, e_2 \in E_{\cmpsortsubscript}, e \in E_{\boxsortsubscript}$ such that $e_1 <^{\mu} x$ and $e_2 <^{\mu} y$, and $e <^{\mu} e_1, e_2$.
Closure under connectivity is then automatic, since $e_1$ and $e_2$ have no incident vertices.
Third, child-less hierarchical edges were forbidden; here a childless
component edge is permitted, and represents the summand $\id_{I}$ of a join---in the original work the treatment of such morphisms was handled by special unlabelled edges that could be introduced and deleted with rewrite rules.
With this formulation the treatment of such morphisms is uniform.
\end{remark}

%% ---------------------------------------------------------------------------
%% Homomorphisms
%% ---------------------------------------------------------------------------

\begin{definition}[Lax, anchored and strict homomorphisms]\label{def:rooted_homomorphism}
A \emph{lax homomorphism} $\phi : \mathcal{F} \to \mathcal{G}$ of rooted
e-hypergraphs is a pair of functions $\phi_{V} : V_{\mathcal{F}} \to
V_{\mathcal{G}}$, $\phi_{E} : E_{\mathcal{F}} \to E_{\mathcal{G}}$ such that
\begin{enumerate}
  \item $\phi_{V}^{*}(s_{\mathcal{F}}(e)) = s_{\mathcal{G}}(\phi_{E}(e))$ and
        $\phi_{V}^{*}(t_{\mathcal{F}}(e)) = t_{\mathcal{G}}(\phi_{E}(e))$;
  \item $l_{\mathcal{F}}(e) = l_{\mathcal{G}}(\phi_{E}(e))$;
  \item if $e <_{\mathcal{F}} x$ then $\phi(e) <_{\mathcal{G}} \phi(x)$.
\end{enumerate}
We call $\phi$ \emph{anchored} if it moreover preserves immediate parents,
\begin{enumerate}
  \item[(3$'$)] if $e <^{\mu}_{\mathcal{F}} x$ then
        $\phi(e) <^{\mu}_{\mathcal{G}} \phi(x)$,
\end{enumerate}
in which case $\phi(\top_{\mathcal{F}})$, a component edge by condition (2), is
called the \emph{anchor} of $\phi$; and \emph{strict} if it is anchored and
$\phi(\top_{\mathcal{F}}) = \top_{\mathcal{G}}$.
Rooted e-hypergraphs with lax homomorphisms form a category
$\catname{EHyp}_{\leq}(\Sigma)$; anchored and strict homomorphisms are closed
under composition and contain the identities, giving wide subcategories
\[
\catname{EHyp}_{\mathrm{s}}(\Sigma)
\;\subseteq\;
\catname{EHyp}_{\bullet}(\Sigma)
\;\subseteq\;
\catname{EHyp}_{\leq}(\Sigma)~.
\]
\end{definition}

\begin{remark}[Comparison with~\cite{tiurin2025equivalencehypergraphsdporewriting}]\label{remark:compare_homomorphism}
In~\cite[Def.~IV.3]{tiurin2025equivalencehypergraphsdporewriting}, a homomorphism of
e-hypergraphs preserves immediate parents \emph{where they exist} and
preserves the consistency relation; parent-less elements are individually
unconstrained, so a homomorphism may distribute two top-level elements into
unrelated positions of the hierarchy.
The three classes above replace this single notion, and no condition is
waived at the root in any of them: since component edges ($e \in E_{\cmpsortsubscript}$) are $0$-ary,
condition (1) holds at $\top_{\mathcal{F}}$ automatically, whatever its
image.
Anchored homomorphisms are \emph{stronger} than the original notion on
former top-level elements: as children of $\top_{\mathcal{F}}$, condition
(3$'$) binds them collectively, placing all of them in the sibling set of a
single component edge---the anchor.
This recovers, as a property of the morphism itself, two side conditions
that~\cite[Def.~V.2]{tiurin2025equivalencehypergraphsdporewriting} imposes on boundary
complements, namely that the images of interface vertices share their
parents and be mutually consistent.
Lax homomorphisms are \emph{weaker}: preserving the ancestor order but not
immediate parenthood, they permit the children of $\top_{\mathcal{F}}$ to be
distributed across distinct components.
This is precisely the behaviour required of internal interface legs of
extended cospans; in~\cite[Def.~IV.5]{tiurin2025equivalencehypergraphsdporewriting} such legs are
homomorphisms only vacuously, their domains being discrete and carrying no
hierarchical structure, whereas here they are morphisms of
$\catname{EHyp}_{\leq}(\Sigma)$ with genuine content.
Finally, preservation of consistency is not a clause of any of the three
classes: for anchored homomorphisms it is a consequence of (3$'$) and
condition (2), since siblings map to siblings of the anchor.
\end{remark}

\begin{definition}[Sub-e-hypergraph]\label{def:sub_ehypergraph}
Let $\mathcal{G}$ be a rooted e-hypergraph.
A \emph{sub-e-hypergraph} of $\mathcal{G}$ is a pair of subsets
$V' \subseteq V_{\mathcal{G}}$, $E' \subseteq E_{\mathcal{G}}$ together with a
distinguished $r \in E'$ such that
\begin{enumerate}
  \item for every $e \in E'$, every vertex occurring in $s_{\mathcal{G}}(e)$ or
        $t_{\mathcal{G}}(e)$ lies in $V'$;
  \item $l_{\mathcal{G}}(r) = \cmpsort$, and every
        $x \in (V' \coprod E') \setminus \{\kappa^{E'}(r)\}$ satisfies
        ${<^{\mu}_{\mathcal{G}}}(x) \in E'$;
  \item every $\boxsort$-edge in $E'$ has at least two immediate children in
        $E'$,
\end{enumerate}
where $\kappa^{E'}$ is the injection of $E'$ into $V' \coprod E'$.
We regard it as the rooted e-hypergraph
$(V', E', s_{\mathcal{G}} \restriction_{E'}, t_{\mathcal{G}} \restriction_{E'},
l_{\mathcal{G}} \restriction_{E'}, r, {<^{\mu}_{\mathcal{G}}} \restriction)$,
and call $r$ its root.
\end{definition}
Put simply, a sub-e-hypergraph is formed by subsets of vertices and edges such that restricted functions and relation $<^{\mu}$ satisfy the e-hypergraph conditions, the latter in particular requires us to pick a distinguished edge $r$ as the root of the sub-e-hypergraph.

% \begin{remark}[Isomorphisms are strict]\label{remark:iso_strict}
% Every isomorphism of $\catname{EHyp}_{\bullet}(\Sigma)$ is strict: if $\phi$
% is invertible and $\phi(e) = \top_{\mathcal{G}}$ for some
% $e \neq \top_{\mathcal{F}}$, then $e$ has an immediate parent, whose image
% under $\phi$ would be a parent of $\top_{\mathcal{G}}$, contradicting
% $[\top_{\mathcal{G}}) = \varnothing$.
% % The analogous statement fails in $\catname{EHyp}_{\leq}(\Sigma)$: a bijective
% % lax homomorphism need not have a lax inverse, as preservation of $<$ does not
% % reflect along bijections.
% % For this reason, isomorphism classes of cospans below are taken with respect
% % to $\catname{EHyp}_{\mathrm{s}}(\Sigma)$; the notion of \enquote{the same
% % e-hypergraph} is therefore unaffected by admitting non-strict morphisms.
% \end{remark}

%% ---------------------------------------------------------------------------
%% Interfaces
%% ---------------------------------------------------------------------------

\begin{definition}[Pointed discrete e-hypergraphs]\label{def:pointed_discrete}
A rooted e-hypergraph is \emph{pointed discrete} if $E = \{\top\}$, and
\emph{ordered} if additionally equipped with a total order on its vertices.
% Given pointed discrete ordered $n$ and $m$, we write $n \sqcup m$ for the
% pushout of $n \leftarrow T \rightarrow m$ in
% $\catname{EHyp}_{\mathrm{s}}(\Sigma)$, which identifies the two roots and
% concatenates the vertex orders, and $n \setminus m$ for the pointed discrete
% e-hypergraph obtained by removing the vertices of $m$ when $m$ is a
% sub-e-hypergraph of $n$.
% An \emph{interface leg} is a lax homomorphism $\phi : n \to \mathcal{G}$ from
% a pointed discrete $n$ such that $\phi(\top_{n}) = \top_{\mathcal{G}}$.
\end{definition}

% \begin{remark}\label{remark:compare_interfaces}
% From a pointed discrete domain, condition (3) of
% \cref{def:rooted_homomorphism} requires only that vertex images be proper
% descendants of the root, which excludes nothing; an interface leg therefore
% amounts to a function on vertices, matching the effectively unconstrained
% interface legs of~\cite[Def.~IV.5]{tiurin2025equivalencehypergraphsdporewriting}.
% An interface leg is strict precisely when all images are immediate children of
% $\top_{\mathcal{G}}$; this replaces the requirement of op.\ cit.\ that the
% images of external interfaces consist of top-level vertices.
% % Note that the feet of cospans must be pointed for composition to be
% % well-defined: a pushout along an unrooted foot would leave the two carriers'
% % roots unidentified, producing two parentless component edges and violating
% % condition (2) of \cref{def:rooted_ehypergraph}.
% % Identifying the roots through the foot replaces the coproduct of carriers
% % used in op.\ cit.\ by a fibered coproduct over $T$---the coproduct of the
% % coslice $T / \catname{EHyp}_{\mathrm{s}}(\Sigma)$---with $T$ in place of
% % the empty e-hypergraph as monoidal unit.
% \end{remark}

% \begin{remark}
% Equivalently, for a pointed discrete $n$ one could consider a functor $D : \mathbb{F} \to \catname{EHyp}(\Sigma)$ as in~\cref{def:hypI} that would map a coproduct in $\mathbb{F}$ to a fibered coproduct in $\catname{EHyp}(\Sigma)$ and map a set $[n]$ to a rooted e-hypergraph with $n$ vertices such that $\top$ is the immediate parent of all vertices.
% \end{remark}

Our aim is to eventually interpret morphisms of $\mathbf{S}_{\Sigma}^{\join}$ and $\mathbf{S}_{\Sigma,\mathcal{E}}^{\join}$ as rooted e-hypergraphs.
For this we need to know which vertices of the e-hypergraph belong to the domain of the morphism and which to the codomain.
This is done by considering cospans of e-hypergraphs with pointed discrete feet, similarly to hypergraphs with interfaces $\catname{HypI}(\Sigma)$ from~\cref{def:hypI}.
Because our rooted e-hypergraphs are hierarchical, we will need a more complicated notion of interface: nested vertices and edges that are immediate children of edges in $E_{\cmpsortsubscript}$ form sub-e-hypergraphs that intuitively correspond to components of a semilattice join; for these sub-e-hypergraphs to correspond to morphisms, we must know what their input and output vertices are.
These cospans are ought to form a symmetric monoidal category.
In the case of $\catname{HypI}(\Sigma)$, the source category for cospans was presheaf-like and hence had all colimits, in particular coproducts that equipped the category with the monoidal structure and pushouts that defined the composition of cospans.
The category of rooted e-hypergraphs $\catname{EHyp}_{\leq}(\Sigma)$ is not presheaf-like: in particular, it does not have all pushouts.

\begin{example}\label{example:pushout_failure}
Consider the span $X \xleftarrow{f} Z \xrightarrow{g} Y$ in $\catname{EHyp}_{\leq}(\Sigma)$ given by the following rooted e-hypergraphs
\[
\adjustbox{max width=0.45\textwidth}{
\tikzfig{./figures/monoidal_egraphs/pushout_failure_example}
}~.
\]
The span maps are given by the dashed arrows.
In particular, $f_{V}(z) = x$, $g_{V}(z) = y$, $f_{E}(\top_{Z}) = \top_{X}$, and $g_{E}(\top_{Z}) = \top_{Y}$.
The pushout of this span does not exist in $\catname{EHyp}_{\leq}(\Sigma)$.
Suppose that there exists a rooted e-hypergraph $P$ and lax
homomorphisms $j_X:X\to P$ and $j_Y:Y\to P$ forming a pushout cocone.
For the square to commute, the images of $x$ and $y$ in $P$ must be
identified; write
\[
v\defeq j_X^V(x)=j_Y^V(y).
\]
Since $b_X<_Xx$ and $b_Y<_Yy$, laxness gives
\[
j_X^E(b_X)<_Pv,
\qquad
j_Y^E(b_Y)<_Pv.
\]
The ancestors of an element in a rooted tree are linearly ordered.
Consequently, there are three possibilities:
\[
j_X^E(b_X)<_Pj_Y^E(b_Y),\qquad
j_Y^E(b_Y)<_Pj_X^E(b_X),\qquad\text{or}\qquad
j_X^E(b_X)=j_Y^E(b_Y).
\]

None of these choices is universal. For example, suppose that
\[
j_X^E(b_X)<_Pj_Y^E(b_Y).
\]
There is another cocone $q_X:X\to Q$, $q_Y:Y\to Q$ in which the order
is reversed:
\[
q_Y^E(b_Y)<_Qq_X^E(b_X).
\]
Such a cocone is obtained by nesting the copy of $b_X$ inside the
component $p_Y$ and identifying the images of $x$ and $y$.

If $P$ were a pushout, there would be a lax homomorphism $u:P\to Q$
such that
\[
j_X\fatsemi u=q_X,
\qquad
j_Y\fatsemi u=q_Y.
\]
Laxness of $u$ would imply
\[
q_X^E(b_X)
=
u_E(j_X^E(b_X))
<
_Q
u_E(j_Y^E(b_Y))
=
q_Y^E(b_Y),
\]
contradicting $q_Y^E(b_Y)<_Qq_X^E(b_X)$.
The case $j_Y^E(b_Y)<_Pj_X^E(b_X)$ is symmetric.

Finally, if $j_X^E(b_X)=j_Y^E(b_Y)$, factorisation through either
cocone above would require the same edge of $P$ to be mapped both to
$q_X^E(b_X)$ and to $q_Y^E(b_Y)$, which are distinct. This is also
impossible. Hence the span has no pushout in
$\catname{EHyp}_{\leq}(\Sigma)$.
\end{example}

Fortunately, we will be interested only in pushouts of spans of a particular form and for these spans the pushout always exists.
We characterise these spans and the pushout construction for them next.

\begin{proposition}\label{prop:pushout_exists}
Consider the following span in $\catname{EHyp}_{\mathrm{\leq}}(\Sigma)$
% https://q.uiver.app/#q=WzAsMyxbMCwwLCJaIl0sWzIsMCwiWCJdLFswLDIsIlkiXSxbMCwxLCJmIl0sWzAsMiwiZyIsMl1d
\[\begin{tikzcd}
	Z && X \\
	\\
	Y
	\arrow["f", from=1-1, to=1-3]
	\arrow["g"', from=1-1, to=3-1]
\end{tikzcd}\]
such that
\begin{enumerate}
\item $Z$ is pointed discrete;
\item $f : Z \to X$ is strict;
\item $g : Z \to Y$ is anchored.
\end{enumerate}
Then the quotient $X \coprod_{Z} Y$ described below is the pushout of this span in $\catname{EHyp}_{\leq}(\Sigma)$ and $\catname{EHyp}_{\bullet}(\Sigma)$.
\[\begin{tikzcd}
	Z && X & \\
	\\
	Y && {X \coprod_{Z} Y} \\
	&&& Q
	\arrow["f", from=1-1, to=1-3]
	\arrow["g"', from=1-1, to=3-1]
	\arrow["{\widetilde{\iota_X}}", from=1-3, to=3-3]
	\arrow["{j_X}", curve={height=-12pt}, from=1-3, to=4-4]
	\arrow["{\widetilde{\iota_Y}}"', from=3-1, to=3-3]
	\arrow["{j_Y}"', curve={height=18pt}, from=3-1, to=4-4]
	\arrow["\lrcorner"{anchor=center, pos=0.125, rotate=180}, draw=none, from=3-3, to=1-1]
	\arrow["{u!}"', from=3-3, to=4-4]
\end{tikzcd}\]
The coprojection $\widetilde{\iota_X} : X \to X \coprod_{Z} Y$ is anchored and $\widetilde{\iota_Y} : Y \to X \coprod_{Z} Y$ is strict.
If $g$ is furthermore strict, then both coprojections are strict and $X \coprod_{Z} Y$ is the pushout in $\catname{EHyp}_{\mathrm{s}}(\Sigma)$.
\end{proposition}
\begin{proof}
Consider the diagram below,
% https://q.uiver.app/#q=WzAsNSxbMCwwLCJaIl0sWzIsMCwiWCJdLFswLDIsIlkiXSxbMiwyLCJYIFxcY29wcm9kX3tafSBZIl0sWzMsMywiUSJdLFswLDEsImYiXSxbMCwyLCJnIiwyXSxbMSwzLCJcXGlvdGFfMSAvXFxzaW0gUiJdLFsyLDMsIlxcaW90YV8yL1xcc2ltIFIiLDJdLFszLDQsInUhIiwyXSxbMSw0LCJqXzEiLDAseyJjdXJ2ZSI6LTJ9XSxbMiw0LCJqXzIiLDIseyJjdXJ2ZSI6M31dLFszLDAsIiIsMix7InN0eWxlIjp7Im5hbWUiOiJjb3JuZXIifX1dXQ==
\[\begin{tikzcd}
	Z && X & \\
	\\
	Y && {X \coprod_{Z} Y} \\
	&&& Q
	\arrow["f", from=1-1, to=1-3]
	\arrow["g"', from=1-1, to=3-1]
	\arrow["{\widetilde{\iota_X}}", from=1-3, to=3-3]
	\arrow["{j_X}", curve={height=-12pt}, from=1-3, to=4-4]
	\arrow["{\widetilde{\iota_Y}}"', from=3-1, to=3-3]
	\arrow["{j_Y}"', curve={height=18pt}, from=3-1, to=4-4]
	\arrow["\lrcorner"{anchor=center, pos=0.125, rotate=180}, draw=none, from=3-3, to=1-1]
	\arrow["{u!}"', from=3-3, to=4-4]
\end{tikzcd}\]
The pushout $X \coprod_{Z} Y$ is computed in two steps.
First, we compute the coproducts of the underlying sets
\[
X \coprod Y = (V_{X} \coprod V_{Y}, E_{X} \coprod E_{Y}, s_{X \coprod Y}, t_{X \coprod Y}, l_{X \coprod Y}, <^{\mu}_{X \coprod Y})~,
\]
where $s_{X \coprod Y} : (E_{X} \coprod E_{Y}) \to (V_{X} \coprod V_{Y})^{*}$ is given by copairing 
\[
[s'_{X}, s'_{Y}] : E_{X} \coprod E_{Y} \to (V_{X} \coprod V_{Y})^{*}
\]
and $s'_{X} : E_{X} \to (V_{X} \coprod V_{Y})^{*}$, $s'_{Y} : E_{Y} \to (V_{X} \coprod V_{Y})^{*}$ defined as 
\begin{gather*}
s'_{X} = s_{X} \fatsemi (\iota_{X}^{V})^{*}\\
s'_{Y} = s_{Y} \fatsemi (\iota_{Y}^{V})^{*}~;
\end{gather*}
$(-)^{*}$ being a monad in particular means it lifts functions $\iota^{V}_{X} : V_{X} \to V_{X} \coprod V_{Y}$ to functions $(\iota_{X}^{V})^{*} : V_{X}^{*} \to (V_{X} \coprod V_{Y})^{*}$.
Similarly for $t_{X \coprod Y}$. The label function $l_{X \coprod Y}$ is given by $[l_{X}, l_{Y}]$. 
The parent relations $<^{\mu}_{Z}$, $<^{\mu}_{X}$ and $<^{\mu}_{Y}$ can be viewed as partial functions, defined on all vertices and edges except for $\top$: that is, we can write $e <^{\mu} x$ meaning $<^{\mu}(x) = e$.
Then, $<^{\mu}_{X \coprod Y} : (V_{X \coprod Y} \coprod E_{X \coprod Y}) \rightharpoonup E_{X \coprod Y}$ is the copairing of partial functions $[\iota^{E}_{X} \circ <^{\mu}_{X}, \iota^{E}_{Y} \circ <^{\mu}_{Y}]$; we freely identify iterated finite coproducts using their canonical associativity and symmetry isomorphisms, i.e., we silently assume $(V_{X \coprod Y} \coprod E_{X \coprod Y}) = (V_{X} \coprod E_{X}) \coprod (V_{Y} \coprod E_{Y})$ without writing down the isomorphism explicitly.
We will also write $\kappa^{V}_{\mathcal{G}} : V_{\mathcal{G}} \to (V_{\mathcal{G}} \coprod E_{\mathcal{G}})$ and $\kappa_{\mathcal{G}}^{E} : E_{\mathcal{G}} \to V_{\mathcal{G}} \coprod  E_{\mathcal{G}}$ for the corresponding injections and in particular we have that
\begin{gather*}
<^{\mu}_{X \coprod Y}(\kappa_{X \coprod Y}^{V}(\iota_{X}^{V}(v))) = \iota_{X}^{E}(<^{\mu}_{X}(\kappa_{X}^{V}(v)))\\
<^{\mu}_{X \coprod Y}(\kappa_{X \coprod Y}^{E}(\iota_{X}^{E}(e))) = \iota_{X}^{E}(<^{\mu}_{X}(\kappa_{X}^{E}(e)))
\end{gather*}
whenever the right-hand side is defined and similarly for $Y$.
Note that $X \coprod Y$ is not a rooted e-hypergraph---it does not have a distinguished $\top$.
In the proof, we will write down all injections explicitly.

Consider the following relations on $V_{X} \coprod V_{Y}$ and $E_{X} \coprod E_{Y}$:
\begin{gather*}
S_{V} = \{\iota^{V}_{X}(v_1) \sim \iota^{V}_{Y}(v_2) \mid \exists v \in V_{Z} : f_{V}(v) = v_1 \wedge g_{V}(v) = v_2\}~,\\
S_{E} = \{\iota^{E}_{X}(e_1) \sim \iota^{E}_{Y}(e_2) \mid f_{E}(\top_{Z}) = e_1 \wedge g_{E}(\top_{Z}) = e_2\}~.
\end{gather*}
and let $R_{V}$ and $R_{E}$ be their reflexive, symmetric and transitive closures.
Note that the content of $S_{E}$ is such because $Z$ is discrete, i.e., $E_{Z} = \{\top_{Z}\}$.
We claim that
\begin{align*}
X \coprod_{Z} Y = \bigl(&V_{X} \coprod V_{Y} / \sim R_{V},\\
 &E_{X} \coprod E_{Y} / \sim R_{E},\\
 &s_{X \coprod Y} / \sim (R_{V}, R_{E}),\\
 &t_{X \coprod Y} / \sim (R_{V}, R_{E}),\\
 &l_{X \coprod Y} / \sim (R_{V}, R_{E}),\\
 &[\iota^{E}_{Y}(\top_{Y})]_{R_{E}},\\
 &<^{\mu}_{X \coprod Y} / \sim (R_{V}, R_{E})\bigr)~.
\end{align*}
is the pushout of the span $X \xleftarrow{f} Z \xrightarrow{g} Y$ in $\catname{EHyp}_{\leq}(\Sigma)$,
where $V_{X} \coprod V_{Y} / \sim R_{V}, E_{X} \coprod E_{Y} / \sim R_{E}$ are the quotient sets and $s_{X \coprod Y} / \sim (R_{V}, R_{E})$, $t_{X \coprod Y} / \sim (R_{V}, R_{E})$ and $<^{\mu}_{X \coprod Y} / \sim (R_{V}, R_{E})$ are the induced functions and partial functions on the quotient sets; we also let $[v]_{R_{V}}$ and $[e]_{R_{E}}$ denote the equivalence classes of $v \in V_{X} \coprod V_{Y}$ and $e \in E_{X} \coprod E_{Y}$ under the relations $R_{V}$ and $R_{E}$ respectively. 

In more detail, let $[-]_{V} : V_{X} \coprod V_{Y} \to V_{X} \coprod V_{Y} / \sim R_{V}$ be the canonical surjection, and similarly for $[-]_{E} : E_{X} \coprod E_{Y} \to E_{X} \coprod E_{Y} / \sim R_{E}$.
Then we define 
\begin{gather}\label{def:quotient_source}
  s_{X \coprod Y} / \sim (R_{V}, R_{E}) : E_{X} \coprod E_{Y} / \sim R_{E} \to (V_{X} \coprod V_{Y} / \sim R_{V})^{*}\\
  [e] \mapsto [(s_{X \coprod Y}(e))]^{*}_{V}~,\notag
\end{gather}
\begin{gather}\label{def:quotient_target}
  t_{X \coprod Y} / \sim (R_{V}, R_{E}) : E_{X} \coprod E_{Y} / \sim R_{E} \to (V_{X} \coprod V_{Y} / \sim R_{V})^{*}\\
  [e] \mapsto [(t_{X \coprod Y}(e))]^{*}_{V}~,\notag
\end{gather}
\begin{gather}\label{def:quotient_label}
  l_{X \coprod Y} / \sim (R_{V}, R_{E}) : E_{X} \coprod E_{Y} / \sim R_{E} \to \Sigma_{1} \cup \{\boxsort, \cmpsort\}\\
  [e] \mapsto l_{X \coprod Y}(e)~,\notag
\end{gather}
\begin{gather}\label{def:quotient_parent}
  <^{\mu}_{X \coprod Y} / \sim (R_{V}, R_{E}) : (V_{X} \coprod V_{Y} / \sim R_{V}) \coprod (E_{X} \coprod E_{Y} / \sim R_{E}) \rightharpoonup E_{X} \coprod E_{Y} / \sim R_{E}\\
  (<^{\mu}_{X \coprod Y} / \sim) \bigl(\kappa_{X \coprod_{Z} Y}^{V}([v])\bigr) = \begin{cases}\notag
    [<^{\mu}_{X \coprod Y}(\kappa_{X \coprod Y}^{V}(v'))]_{E} & 
    \begin{aligned}
    &\text{if } \exists\; v' \in [v]_{V} \text{ such that }\\
    &<^{\mu}_{X \coprod Y}(\kappa_{X \coprod Y}^{V}(v')) \text{ is defined}~,
    \end{aligned}
    \\[2mm]
    \text{undefined} & \text{otherwise}
  \end{cases}\\ 
  (<^{\mu}_{X \coprod Y} / \sim) \bigl(\kappa_{X \coprod_{Z} Y}^{E}([e])\bigr) = \begin{cases}
    [<^{\mu}_{X \coprod Y}(\kappa_{X \coprod Y}^{E}(e'))]_{E} & \begin{aligned}
    &\text{if } \exists\; e' \in [e]_{E} \text{ such that }\\
    &<^{\mu}_{X \coprod Y}(\kappa_{X \coprod Y}^{E}(e')) \text{ is defined}~,
    \end{aligned}
    \\[2mm]
    \text{undefined} & \text{otherwise}
  \end{cases}\notag 
\end{gather}
Moving forward, we will refer to either $\sim (R_{V}, R_{E})$, $\sim R_{V}$ or $\sim R_{E}$ as $\sim$.
We will also write $X \coprod_{Z} Y$ as 
\[ 
X \coprod_{Z} Y = (V_{X \coprod_{Z} Y}, E_{X \coprod_{Z} Y}, s_{X \coprod_{Z} Y}, t_{X \coprod_{Z} Y}, l_{X \coprod_{Z} Y},  \top_{X \coprod_{Z} Y}, <^{\mu}_{X \coprod_{Z} Y}) 
\]
for brevity, and use that to refer to its constituent components.
Let us check if the functions above are well-defined.

For $s_{X \coprod Y}/\sim$ we need to show that if $[e] = [e']$ then $[(s_{X \coprod Y}(e))]^{*}_{V} = [(s_{X \coprod Y}(e'))]^{*}_{V}$.
By definition of $R_{E}$, $[e] = [e']$ for $e \not = e'$ implies that they are in the image of $\top_{Z}$.
Without loss of generality, let $e = \iota^{E}_{X}(f_{E}(\top_{Z}))$ and $e' = \iota^{E}_{Y}(g_{E}(\top_{Z}))$.
By definition, $s_{X \coprod Y}(\iota^{E}_{X}(f_{E}(\top_{Z}))) = (\iota^{V}_{X})^{*}(s_{X}(f_{E}(\top_{Z})))$ and because $f = (f_{V}, f_{E})$ is a homomorphism, $s_{X}(f_{E}(\top_{Z})) = f_{V}^{*}(s_{Z}(\top_{Z})) = \varepsilon$ and hence $s_{X \coprod Y}(\iota^{E}_{X}(f_{E}(\top_{Z}))) = \varepsilon$.
Similarly, $s_{X \coprod Y}(\iota^{E}_{Y}(g_{E}(\top_{Z}))) = (\iota^{V}_{Y})^{*}(s_{Y}(g_{E}(\top_{Z}))) = (\iota^{V}_{Y})^{*}(g_{V}^{*}(s_{Z}(\top_{Z}))) = \varepsilon$.
The same argument applies to $t_{X \coprod Y}/\sim$.

For the label function we have that
\[
l_{X}(f_{E}(\top_{Z})) = l_{Z}(\top_{Z}) = \cmpsort = l_{Y}(g_{E}(\top_{Z}))~.
\] and by definition 
\[
l_{X \coprod Y}(\iota^{E}_{X}(f_{E}(\top_{Z}))) = l_{X}(f_{E}(\top_{Z})) = \cmpsort~.
\]
Similarly, $l_{X \coprod Y}(\iota^{E}_{Y}(g_{E}(\top_{Z}))) = l_{Y}(g_{E}(\top_{Z})) = \cmpsort$.
Well-definedness of $<^{\mu}_{X \coprod_{Z} Y}$ is given by~\cref{lemma:immediate_predecessor_defined}.
Then, $X \coprod_{Z} Y$ is a rooted e-hypergraph by~\cref{lemma:quotient_is_e-hypergraph}.
Maps $\widetilde{\iota_{X}}$ and $\widetilde{\iota_{Y}}$ are lax homomorphisms by~\cref{lemma:pushout_injections_are_homomorphisms} and moreover $\widetilde{\iota_{X}}$ is anchored and $\widetilde{\iota_{Y}}$ is strict by the same lemma and finally both homomorphisms are strict if $g$ is strict.
Finally, $(X \coprod_{Z} Y, \widetilde{\iota_{X}}, \widetilde{\iota_{Y}})$ has the universal property of the pushout, that is for any other rooted e-hypergraph $Q$ with candidate maps $j_{X}$ and $j_{Y}$, there is a unique lax homomorphism $u$ such that the appropriate triangles commute.
This $u$ is anchored when both $j_{X}$ and $j_{Y}$ are anchored, and strict when $j_{Y}$ is strict and $j_{X}$ is at least anchored.
This is given by~\cref{lemma:pushout_universal}.
\end{proof}


\begin{lemma}\label{lemma:immediate_predecessor_defined}
The partial function $<^{\mu}_{X \coprod Y} / \sim : (V_{X} \coprod V_{Y} / \sim) \coprod (E_{X} \coprod E_{Y} / \sim) \rightharpoonup E_{X} \coprod E_{Y} / \sim$ is well-defined and does not depend on the choice of representative in the equivalence class.
\end{lemma}
\begin{proof}
Recall the definition of $<^{\mu}_{X \coprod Y} / \sim$:
\begin{gather*}
  (<^{\mu}_{X \coprod Y} / \sim) \bigl(\kappa_{X \coprod_{Z} Y / \sim}^{V}([v])\bigr) = \begin{cases}
    [<^{\mu}_{X \coprod Y}(\kappa_{X \coprod Y}^{V}(v'))]_{E} & 
    \begin{aligned}
    &\text{if } \exists\; v' \in [v]_{V} \text{ such that }\\
    &<^{\mu}_{X \coprod Y}(\kappa_{X \coprod Y}^{V}(v')) \text{ is defined}~,
    \end{aligned}
    \\[2mm]
    \text{undefined} & \text{otherwise}
  \end{cases}\\ 
  (<^{\mu}_{X \coprod Y} / \sim) \bigl(\kappa_{X \coprod_{Z} Y / \sim}^{E}([e])\bigr) = \begin{cases}
    [<^{\mu}_{X \coprod Y}(\kappa_{X \coprod Y}^{E}(e'))]_{E} & \begin{aligned}
    &\text{if } \exists\; e' \in [e]_{E} \text{ such that }\\
    &<^{\mu}_{X \coprod Y}(\kappa_{X \coprod Y}^{E}(e')) \text{ is defined}~,
    \end{aligned}
    \\[2mm]
    \text{undefined} & \text{otherwise}
  \end{cases}~. 
\end{gather*}
We need to show that if $[v_1]_{V} = [v_2]_{V}$ then 
\[
(<^{\mu}_{X \coprod Y} / \sim)\bigl(\kappa_{X \coprod_{Z} Y / \sim}^{V}([v_1])\bigr) =  (<^{\mu}_{X \coprod Y} / \sim)\bigl(\kappa_{X \coprod_{Z} Y / \sim}^{V}([v_2])\bigr)
\] and similarly for edges $[e_1]_{E} = [e_2]_{E}$.
We also need to show that if there exist $u_1 \in [v]_{V}$ and $u_2 \in [v]_{V}$ such that $<^{\mu}_{X \coprod Y}(\kappa_{X \coprod Y}^{V}(u_1))$ and $<^{\mu}_{X \coprod Y}(\kappa_{X \coprod Y}^{V}(u_2))$ are defined, then $[<^{\mu}_{X \coprod Y}(\kappa_{X \coprod Y}^{V}(u_1))]_{E} = [<^{\mu}_{X \coprod Y}(\kappa_{X \coprod Y}^{V}(u_2))]_{E}$ and similarly for edges.
This latter property in fact implies the former, so we will only show the latter.

By definition of $R_{V}$, $u_1 R_{V} u_2$ implies there is a sequence of vertices $u_1 = w_1 \mathcal{S} w_2 \ldots \mathcal{S} w_{n-1} \mathcal{S} w_n = u_2$, where each $w_i \mathcal{S} w_{i+1}$ is either $w_i S_{V} w_{i+1}$, $w_{i+1} S_{V} w_{i}$, or $w_{i} = w_{i+1}$.
We will do induction on the length of this sequence.
The base case is when the length is 2: in the case when $u_1 = w_1 = w_2 = u_2$ there is nothing to prove, so we consider the case $u_1 S_{V} u_2$.
By definition of $S_{V}$, there exists a vertex $z \in V_{Z}$ such that $\iota_{X}^{V}(f_{V}(z)) = u_1$ and $\iota_{Y}^{V}(g_{V}(z)) = u_2$.
By assumption, both $<^{\mu}_{X \coprod Y}(\kappa_{X \coprod Y}^{V}(u_1))$ and $<^{\mu}_{X \coprod Y}(\kappa_{X \coprod Y}^{V}(u_2))$ are defined; let their values be $e_1$ and $e_2$ respectively.
We need to check that $e_1 \sim e_2$.
We have that
\begin{align*}
<^{\mu}_{X \coprod Y}\bigl(\kappa_{X \coprod Y}^{V}( u_1 )\bigr) &= <_{X \coprod Y}^{\mu}(\kappa_{X \coprod Y}^{V}(\iota_{X}^{V}(f_{V}(z))))\\
&= \iota_{X}^{E}(<_{X}^{\mu}(\kappa_{X}^{V}(f_{V}(z))))\\ 
&= \iota_{X}^{E}(f_{E}(<_{Z}^{\mu}(\kappa_{Z}^{V}(z))))\\
&= \iota_{X}^{E}(f_{E}(\top_{Z}))~,
\end{align*}
because $f$ is strict and $Z$ is pointed discrete.
In more detail, $\top_{Z} <_{Z}^{\mu} z$ and because $f$ is at least anchored, $f_E(\top_{Z}) <_{X}^{\mu} (f_{V}(z))$ and since the immediate predecessor is unique, we have that $<^{\mu}_{X}(\kappa_{X}^{V}(f_{V}(z))) = f_{E}(\top_{Z})$.
We also have that
\begin{align*}
<^{\mu}_{X \coprod Y}\bigl(\kappa_{X \coprod Y}^{V}(u_2)\bigr) &= <_{X \coprod Y}^{\mu}(\kappa_{X \coprod Y}^{V}(\iota_{Y}^{V}(g_{V}(z))))\\
&= \iota_{Y}^{E}(<_{Y}^{\mu}(\kappa_{Y}^{V}(g_{V}(z))))\\ 
&= \iota_{Y}^{E}(g_{E}(<_{Z}^{\mu}(\kappa_{Z}^{V}(z))))\\
&= \iota_{Y}^{E}(g_{E}(\top_{Z}))~,
\end{align*}
because $g$ is anchored and applying uniqueness as above.
Therefore,
\begin{gather*}
<^{\mu}_{X \coprod Y}\bigl(\kappa_{X \coprod Y}^{V}( u_1 )\bigr) = \iota_{X}^{E}(f_{E}(\top_{Z}))\\
<^{\mu}_{X \coprod Y}\bigl(\kappa_{X \coprod Y}^{V}( u_2 )\bigr) = \iota_{Y}^{E}(g_{E}(\top_{Z}))~.
\end{gather*}
These two edges are related by the definition of $S_{E}$, and hence their $R_{E}$-classes agree.
The case $u_2 S_{V} u_1$ is symmetric.
Say the sequence has length $n > 2$, that is $u_1 = w_1 \mathcal{S} w_2 \mathcal{S} \ldots \mathcal{S} w_{n-1} \mathcal{S} w_n = u_2$.
By the induction hypothesis, we have that $[<^{\mu}_{X \coprod Y}(\kappa_{X \coprod Y}^{V}(w_1))]_{E} = [<^{\mu}_{X \coprod Y}(\kappa_{X \coprod Y}^{V}(w_{n-1}))]_{E}$, because each $w_i$ is a vertex and hence has an immediate predecessor. 
Then by applying the above argument we can derive that $[<^{\mu}_{X \coprod Y}(\kappa_{X \coprod Y}^{V}(w_{n-1}))]_{E} = [<^{\mu}_{X \coprod Y}(\kappa_{X \coprod Y}^{V}(w_n))]_{E}$ and by transitivity we have 
\[
[<^{\mu}_{X \coprod Y}(\kappa_{X \coprod Y}^{V}(u_1))]_{E} = [<^{\mu}_{X \coprod Y}(\kappa_{X \coprod Y}^{V}(u_2))]_{E}~.
\]
Let us now consider the case of edges.
By definition of $R_{E}$, the equivalence class $[e]_{E}$ has at most two elements.
Therefore, $e_1 R_{E} e_2$ implies there either $e_1 = e_2$, $e_1 S_{E} e_2$, or $e_2 S_{E} e_1$.
In the first case there is nothing to prove, so we consider the second case.
By definition of $S_{E}$, without loss of generality, $e_1 = \iota_{X}^{E}(f_{E}(\top_{Z}))$ and $e_2 = \iota_{Y}^{E}(g_{E}(\top_{Z}))$.
By assumption, both $<^{\mu}_{X \coprod Y}(\kappa_{X \coprod Y}^{E}(e_1))$ and $<^{\mu}_{X \coprod Y}(\kappa_{X \coprod Y}^{E}(e_2))$ are defined.
We then have that
\begin{align*}
<^{\mu}_{X \coprod Y}\bigl(\kappa_{X \coprod Y}^{E}(e_1)\bigr) &= <_{X \coprod Y}^{\mu}(\kappa_{X \coprod Y}^{E}(\iota_{X}^{E}(f_{E}(\top_{Z}))))\\
&= \iota_{X}^{E}(<_{X}^{\mu}(\kappa_{X}^{E}(f_{E}(\top_{Z}))))\\
&= \iota_{X}^{E}(<_{X}^{\mu}(\kappa_{X}^{E}(\top_{X})))
\end{align*}
the latter expression is undefined, because $\top_{X}$ has no predecessor and $f$ is strict.
This is a contradiction, and therefore in each equivalence class $[e]_{E}$ there is at most one edge $e$ such that $<^{\mu}_{X \coprod Y}(\kappa_{X \coprod Y}^{E}(e))$ is defined.
The third case is symmetric.
\end{proof}

\begin{lemma}\label{lemma:immediate_predecessor_undefined}
The partial function $<^{\mu}_{X \coprod_{Z} Y}$ is undefined exactly at $\kappa_{X \coprod_{Z} Y}^{E}([\iota_{Y}^{E}(\top_{Y})]_{E})$.
\end{lemma}
\begin{proof}
  It is sufficient first to check that $<^{\mu}_{X \coprod Y}(\kappa_{X \coprod Y}^{E}(e))$ is undefined for all $e R_{E} \iota_{Y}^{E}(\top_{Y})$.
  By definition, $<^{\mu}_{X \coprod Y}(\kappa_{X \coprod Y}^{E}(\iota_{Y}^{E}(\top_{Y}))) = \iota_{Y}^{E}(<^{\mu}_{Y}(\kappa_{Y}^{E}(\top_{Y})))$, which is undefined.
  $e R_{E} \iota_{Y}^{E}(\top_{Y})$ implies that either $e = \iota_{Y}^{E}(\top_{Y})$, in which case the immediate predecessor of $e$ is undefined, or $e S_{E} \iota_{Y}^{E}(\top_{Y})$.
  In the latter case, by definition of $S_{E}$, $e = \iota_{X}^{E}(f_{E}(\top_{Z}))$ and $\iota_{Y}^{E}(\top_{Y}) = \iota_{Y}^{E}(g_{E}(\top_{Z}))$.
  Because $f$ is strict, $f_{E}(\top_{Z}) = \top_{X}$ and hence $<^{\mu}_{X}(\kappa_{X}^{E}(f_{E}(\top_{Z})))$ is undefined.

  Let us now check that $<^{\mu}_{X \coprod_{Z} Y}$ is defined at all other equivalence classes.
  First, since every vertex in $V_{X} \coprod V_{Y}$ has an immediate predecessor, $<^{\mu}_{X \coprod Y} / \sim$ is defined at all equivalence classes of vertices.
  For edges, let $[e]_{E} \not = [\iota_{Y}^{E}(\top_{Y})]_{E}$.
  If $[e]_{E}$ contains a non-root edge of $X$ or $Y$, then that representative has the immediate predecessor.
  Thus, the only remaining possibility is that $[e]_{E} = [\iota_{X}^{E}(\top_{X})]_{E}$.
  Since $f$ is strict, we have that $\top_{X} = f_{E}(\top_{Z})$ and hence $\iota_{X}^{E}(\top_{X}) S_{E} \iota_{Y}^{E}(g_{E}(\top_{Z}))$.
  By assumption, $g_{E}(\top_{Z}) \not = \top_{Y}$ and hence $\iota_{Y}^{E}(g_{E}(\top_{Z}))$ has the immediate predecessor and so does $[e]_{E}$.
\end{proof}


\begin{definition}[Depth]\label{def:depth}
Let $\mathcal{G}$ be a rooted e-hypergraph.
By condition (2) of~\cref{def:rooted_ehypergraph}, iterating $<^{\mu}_{\mathcal{G}}$
from any $x \in V_{\mathcal{G}} \coprod E_{\mathcal{G}}$ reaches $\top_{\mathcal{G}}$
in finitely many steps, so we may define the \emph{depth}
$d_{\mathcal{G}} : V_{\mathcal{G}} \coprod E_{\mathcal{G}} \to \mathbb{N}$ by
\[
d_{\mathcal{G}}(\kappa^{E}_{\mathcal{G}}(\top_{\mathcal{G}})) = 0~,
\qquad
d_{\mathcal{G}}(x) = d_{\mathcal{G}}\bigl(\kappa^{E}_{\mathcal{G}}(<^{\mu}_{\mathcal{G}}(x))\bigr) + 1
\quad\text{for } x \neq \kappa^{E}_{\mathcal{G}}(\top_{\mathcal{G}})~.
\]
\end{definition}

\begin{lemma}[Depth of the quotient]\label{lemma:pushout_depth}
Let $\delta \defeq d_{Y}(\kappa^{E}_{Y}(g_{E}(\top_{Z})))$ be the depth of the
anchor of $g$.
Then there is a function
\[
\mathfrak{D}:
(V_{X \coprod_{Z} Y}) \coprod (E_{X \coprod_{Z} Y}) \to \mathbb{N}
\]
determined by
\begin{gather*}
\mathfrak{D}\bigl(\kappa^{V}_{X \coprod_{Z} Y}([\iota^{V}_{X}(v)]_{V})\bigr) = \delta + d_{X}(\kappa^{V}_{X}(v))~,
\qquad
\mathfrak{D}\bigl(\kappa^{E}_{X \coprod_{Z} Y}([\iota^{E}_{X}(e)]_{E})\bigr) = \delta + d_{X}(\kappa^{E}_{X}(e))~,\\
\mathfrak{D}\bigl(\kappa^{V}_{X \coprod_{Z} Y}([\iota^{V}_{Y}(v)]_{V})\bigr) = d_{Y}(\kappa^{V}_{Y}(v))~,
\qquad
\mathfrak{D}\bigl(\kappa^{E}_{X \coprod_{Z} Y}([\iota^{E}_{Y}(e)]_{E})\bigr) = d_{Y}(\kappa^{E}_{Y}(e))~,
\end{gather*}
and moreover, for every $\xi \in (V_{X \coprod_{Z} Y}) \coprod (E_{X \coprod_{Z} Y})$
at which $<^{\mu}_{X \coprod_{Z} Y}$ is defined,
\[
\mathfrak{D}(\xi)
= \mathfrak{D}\bigl(\kappa^{E}_{X \coprod_{Z} Y}(<^{\mu}_{X \coprod_{Z} Y}(\xi))\bigr) + 1~,
\]
and for $\xi = \kappa^{E}_{X \coprod_{Z} Y}([\iota^{E}_{Y}(\top_{Y})]_{E})$,
\[
\mathfrak{D}(\xi) = 0~.
\]
\end{lemma}
\begin{proof}
We need to show that the function $\mathfrak{D}$ is well-defined on the quotient.
For vertices, pick $v_1 R_{V} v_2$; we need to show that $\mathfrak{D}(\kappa^{V}_{X \coprod_{Z} Y}([v_1]_{V})) = \mathfrak{D}(\kappa^{V}_{X \coprod_{Z} Y}([v_2]_{V}))$.
As before, it suffices to consider the case $v_1 S_{V} v_2$, the other cases follow by symmetry and transitivity.
By definition of $S_{V}$, there exists a vertex $z \in V_{Z}$ such that, without loss of generality, $v_1 = \iota^{V}_{X}(f_{V}(z))$ and $v_2 = \iota^{V}_{Y}(g_{V}(z))$.
We then need to check that $\delta + d_{X}(\kappa^{V}_{X}(f_{V}(z))) = d_{Y}(\kappa^{V}_{Y}(g_{V}(z)))$.
Since $Z$ is pointed discrete, $\top_{Z} <^{\mu}_{Z} z$, and since $f$ is strict, 
\begin{align*}
  d_{X}(\kappa^{V}_{X}(f_{V}(z))) &= d_{X}(\kappa^{E}_{X}(<^{\mu}_{X}(\kappa^{V}_{X}(f_{V}(z))))) + 1\\
  &= d_{X}(\kappa^{E}_{X}(f_{E}(\top_{Z}))) + 1\\
  &= d_{X}(\kappa^{E}_{X}(\top_{X})) + 1\\
  &= 1~.
\end{align*}
Then, recalling that $\delta = d_{Y}(\kappa^{E}_{Y}(g_{E}(\top_{Z})))$ and that $g$ is anchored, we have $g_{E}(\top_{Z}) <^{\mu}_{Y} g_{V}(z)$ and hence 
\begin{align*}
d_{Y}(\kappa^{V}_{Y}(g_{V}(z))) &= d_{Y}(\kappa^{E}_{Y}(g_{E}(\top_{Z}))) + 1\\
&= \delta + 1~.
\end{align*}
For edges pick $e_1 R_{E} e_2$; we need to show that $\mathfrak{D}(\kappa^{E}_{X \coprod_{Z} Y}([e_1]_{E})) = \mathfrak{D}(\kappa^{E}_{X \coprod_{Z} Y}([e_2]_{E}))$.
Consider the case $e_1 S_{E} e_2$.
By definition, $e_1 = \iota^{E}_{X}(f_{E}(\top_{Z}))$ and $e_2 = \iota^{E}_{Y}(g_{E}(\top_{Z}))$.
Then $\mathfrak{D}(\kappa^{E}_{X \coprod_{Z} Y}([e_1]_{E})) = \delta + d_{X}(\kappa^{E}_{X}(f_{E}(\top_{Z}))) = \delta + 0 = \delta$ and $\mathfrak{D}(\kappa^{E}_{X \coprod_{Z} Y}([e_2]_{E})) = d_{Y}(\kappa^{E}_{Y}(g_{E}(\top_{Z}))) = \delta$.

We already established that $\mathfrak{D}(\kappa_{X \coprod_{Z} Y}^{E}([\iota^{E}_{Y}(\top_{Y})]_{E})) = 0$, so it remains to check that
for every $\xi \in (V_{X \coprod_{Z} Y}) \coprod (E_{X \coprod_{Z} Y})$ at which $<^{\mu}_{X \coprod_{Z} Y}$ is defined, we have 
$\mathfrak{D}(\xi) = \mathfrak{D}(\kappa^{E}_{X \coprod_{Z} Y}(<^{\mu}_{X \coprod_{Z} Y}(\xi))) + 1$.
Let $\xi = \kappa^{V}_{X \coprod_{Z} Y}([v]_{V})$ and suppose $[v]_{V}$ has a representative $\iota_{X}^{V}(v')$.
We then have that 
\[
\mathfrak{D}(\xi) = \delta + d_{X}(\kappa^{V}_{X}(v'))~.
\]
On the other hand, $<^{\mu}_{X \coprod_{Z} Y}(\xi) = [<^{\mu}_{X \coprod Y}(\kappa_{X \coprod Y}^{V}(\iota_{X}^{V}(v')))]_{E} = [\iota^{E}_{X}(<^{\mu}_{X}(\kappa_{X}^{V}(v')))]_{E}$.
We then have that 
\begin{align*}
  \mathfrak{D}(\kappa^{E}_{X \coprod_{Z} Y}(<^{\mu}_{X \coprod_{Z} Y}(\xi))) &= \mathfrak{D}(\kappa_{X \coprod_{Z} Y}^{E}([\iota^{E}_{X}(<^{\mu}_{X}(\kappa_{X}^{V}(v')))]_{E}))\\
  &= \delta + d_{X}(\kappa^{E}_{X}(<^{\mu}_{X}(\kappa_{X}^{V}(v'))))\\
  &= \delta + d_{X}(\kappa^{V}_{X}(v')) - 1\\
  &= \mathfrak{D}(\xi) -1~.
\end{align*}
Therefore, we established that $\mathfrak{D}(\xi) = \mathfrak{D}(\kappa^{E}_{X \coprod_{Z} Y}(<^{\mu}_{X \coprod_{Z} Y}(\xi))) + 1$.
The cases in the $Y$-summand are symmetric.
The case for edges follows by noting that the definition of $<^{\mu}_{X \coprod_{Z} Y}$ on edges implies that for an edge $[e]_{E}$ for which the immediate predecessor is defined, there exists a representative $e' \in E_{X \coprod Y}$ such that $<^{\mu}_{X \coprod Y}(\kappa_{X \coprod Y}^{E}(e'))$ is defined and hence the same argument applies.
\end{proof}

\begin{lemma}\label{lemma:quotient_is_e-hypergraph}
The quotient $X \coprod_Z Y$ is a rooted e-hypergraph.
\end{lemma}
\begin{proof}
The vertex and edge sets of $X \coprod_Z Y$ are finite, since they are
quotients of finite coproducts.
The source, target, label and immediate-parent maps are well-defined by
the preceding results.

It remains to check conditions (1)--(5) of
\cref{def:rooted_ehypergraph}.
\begin{enumerate}
  \item We first check the typing condition.

  Let $[e]_E$ be a labelled edge of $X \coprod_Z Y$ such that
  \[
  l_{X \coprod_Z Y}([e]_E) : m \to n.
  \]
  Choose a representative $e \in E_X \coprod E_Y$.
  Without loss of generality, suppose that
  \[
  e = \iota_X^E(e')
  \]
  for some $e' \in E_X$.
  Then
  \[
  l_{X \coprod_Z Y}([e]_E)
  =
  l_{X \coprod Y}(\iota_X^E(e'))
  =
  l_X(e').
  \]
  Hence $l_X(e') : m \to n$.

  By definition of the quotient source map,
  \begin{align*}
  s_{X \coprod_Z Y}([e]_E)
  &=
  ([-]_V)^*
  \bigl(s_{X \coprod Y}(\iota_X^E(e'))\bigr)\\
  &=
  ([-]_V)^*
  \bigl((\iota_X^V)^*(s_X(e'))\bigr).
  \end{align*}
  Both $(\iota_X^V)^*$ and $([-]_V)^*$ preserve the length of
  words. Therefore,
  \[
  \bigl|s_{X \coprod_Z Y}([e]_E)\bigr|
  =
  |s_X(e')|
  =
  m.
  \]
  Similarly,
  \[
  \bigl|t_{X \coprod_Z Y}([e]_E)\bigr|
  =
  |t_X(e')|
  =
  n.
  \]
  The case in which $e$ lies in the $Y$-summand is symmetric.

  Now let $[e]_E$ be a component edge of $X \coprod_Z Y$.
  Choose a representative $e \in E_X \coprod E_Y$ and, without loss
  of generality, suppose that $e = \iota_X^E(e')$.
  Since the quotient label map preserves labels,
  \[
  l_X(e') = \cmpsort.
  \]
  Hence, by the typing condition for $X$,
  \[
  s_X(e') = t_X(e') = \varepsilon.
  \]
  It follows that
  \begin{align*}
  s_{X \coprod_Z Y}([e]_E)
  &=
  ([-]_V)^*
  \bigl((\iota_X^V)^*(s_X(e'))\bigr)
  =
  \varepsilon,\\
  t_{X \coprod_Z Y}([e]_E)
  &=
  ([-]_V)^*
  \bigl((\iota_X^V)^*(t_X(e'))\bigr)
  =
  \varepsilon.
  \end{align*}
  Again, the case in which $e$ lies in the $Y$-summand is symmetric.
\item The root of $X \coprod_Z Y$ is
$[\iota_Y^E(\top_Y)]_E$, and
\[
l_{X \coprod_Z Y}([\iota_Y^E(\top_Y)]_E)
= l_Y(\top_Y)
= \cmpsort.
\]
By~\cref{lemma:immediate_predecessor_defined},
$<^\mu_{X \coprod_Z Y}$ is a partial function defined at every
vertex and edge except the root. Hence every non-root element has
exactly one immediate parent, while the root has none.

Moreover, by~\cref{lemma:pushout_depth}, if
$e <^\mu_{X \coprod_Z Y} \xi$, then
\[
\mathfrak D(\xi)
=
\mathfrak D(\kappa^E_{X\coprod_Z Y}(e))+1.
\]
Thus depth strictly decreases whenever the immediate-parent function
is iterated, so there can be no cycle. Iteration must terminate at an
element on which the parent function is undefined, and the only such
element is the root. Therefore the transitive closure of
$<^\mu_{X\coprod_Z Y}$ is a strict partial order and
$(V_{X\coprod_ZY}\coprod E_{X\coprod_ZY},<^\mu_{X\coprod_ZY})$
is a tree rooted at $[\iota_Y^E(\top_Y)]_E$.
    \item We check alternation; recall that $<^{\mu}_{X \coprod_{Z} Y}$ is defined at
  every element other than the root by~\cref{lemma:immediate_predecessor_defined}.
  We need to check that for any edge
  $[e]_{E} \in (E_{X \coprod_{Z} Y})_{\Sigma} \cup (E_{X \coprod_{Z} Y})_{\boxsortsubscript}$,
  \[
  l_{X \coprod_{Z} Y}\bigl(<^{\mu}_{X \coprod_{Z} Y}(\kappa^{E}_{X \coprod_{Z} Y}([e]_{E}))\bigr) = \cmpsort~;
  \]
  that for any $[v]_{V} \in V_{X \coprod_{Z} Y}$,
  \[
  l_{X \coprod_{Z} Y}\bigl(<^{\mu}_{X \coprod_{Z} Y}(\kappa^{V}_{X \coprod_{Z} Y}([v]_{V}))\bigr) = \cmpsort~;
  \]
  and that for any $[e]_{E} \in (E_{X \coprod_{Z} Y})_{\cmpsortsubscript}$ with
  $[e]_{E} \neq \top_{X \coprod_{Z} Y}$,
  \[
  l_{X \coprod_{Z} Y}\bigl(<^{\mu}_{X \coprod_{Z} Y}(\kappa^{E}_{X \coprod_{Z} Y}([e]_{E}))\bigr) = \boxsort~.
  \]
  All three statements follow from the observation that both the sort of an
  element of $X \coprod_{Z} Y$ and its immediate parent are computed from a
  representative lying in a single summand.
  First, consider the case of vertices.

  Let $\xi = \kappa^{V}_{X \coprod_{Z} Y}([v]_{V})$.
  By~\cref{lemma:immediate_predecessor_defined} the immediate parent of $\xi$ is
  defined, and by the construction of $<^{\mu}_{X \coprod_{Z} Y}$ there is a
  representative $v' \in [v]_{V}$ at which $<^{\mu}_{X \coprod Y}$ is defined and
  such that
  \[
  <^{\mu}_{X \coprod_{Z} Y}(\xi) = [<^{\mu}_{X \coprod Y}(\kappa^{V}_{X \coprod Y}(v'))]_{E}~.
  \]
  Without loss of generality $v' = \iota^{V}_{X}(v'')$ for some $v'' \in V_{X}$,
  the case of the $Y$-summand being symmetric.
  Then
  \[
  <^{\mu}_{X \coprod Y}(\kappa^{V}_{X \coprod Y}(\iota^{V}_{X}(v'')))
  = \iota^{E}_{X}\bigl(<^{\mu}_{X}(\kappa^{V}_{X}(v''))\bigr)~,
  \]
  and hence, by the definition of the quotient label map,
  \begin{align*}
  l_{X \coprod_{Z} Y}\bigl(<^{\mu}_{X \coprod_{Z} Y}(\xi)\bigr)
  &= l_{X \coprod_{Z} Y}\bigl([\iota^{E}_{X}(<^{\mu}_{X}(\kappa^{V}_{X}(v'')))]_{E}\bigr)\\
  &= l_{X \coprod Y}\bigl(\iota^{E}_{X}(<^{\mu}_{X}(\kappa^{V}_{X}(v'')))\bigr)
   = l_{X}\bigl(<^{\mu}_{X}(\kappa^{V}_{X}(v''))\bigr)~.
  \end{align*}
  Since $X$ is a rooted e-hypergraph and $v''$ is a vertex, condition (3)
  of~\cref{def:rooted_ehypergraph} gives
  $l_{X}(<^{\mu}_{X}(\kappa^{V}_{X}(v''))) = \cmpsort$, which is the second of
  the three statements.

  Now let $\xi = \kappa^{E}_{X \coprod_{Z} Y}([e]_{E})$ with
  $[e]_{E} \neq \top_{X \coprod_{Z} Y}$.
  As above there is a representative $e' \in [e]_{E}$ at which
  $<^{\mu}_{X \coprod Y}$ is defined, and, without loss of generality,
  $e' = \iota^{E}_{X}(e'')$ for some $e'' \in E_{X}$ with $e'' \neq \top_{X}$.
  The same computation, with $\kappa^{E}$ in place of $\kappa^{V}$, yields
  \[
  l_{X \coprod_{Z} Y}\bigl(<^{\mu}_{X \coprod_{Z} Y}(\xi)\bigr)
  = l_{X}\bigl(<^{\mu}_{X}(\kappa^{E}_{X}(e''))\bigr)~,
  \]
  while
  \[
  l_{X \coprod_{Z} Y}([e]_{E}) = l_{X \coprod Y}(\iota^{E}_{X}(e'')) = l_{X}(e'')~,
  \]
  so $[e]_{E}$ and $e''$ have the same sort.
  Applying condition (3) of~\cref{def:rooted_ehypergraph} in $X$ therefore gives
  the remaining two statements: if $[e]_{E}$ is a labelled or a box edge then so
  is $e''$, therefore $l_{X}(<^{\mu}_{X}(\kappa^{E}_{X}(e''))) = \cmpsort$; and if
  $[e]_{E}$ is a component edge then so is $e''$, therefore
  $l_{X}(<^{\mu}_{X}(\kappa^{E}_{X}(e''))) = \boxsort$.

\item We now check incidence. For every
\[
[e]_{E}
\in
(E_{X \coprod_Z Y})_{\Sigma}
\cup
(E_{X \coprod_Z Y})_{\boxsortsubscript}
\]
and every $[v]_{V} \in V_{X \coprod_Z Y}$, if $[v]_{V}$ occurs in
\[
s_{X \coprod_Z Y}([e]_{E})
\qquad\text{or}\qquad
t_{X \coprod_Z Y}([e]_{E}),
\]
we must show that
\[
<^{\mu}_{X \coprod_Z Y}
\bigl(\kappa^{E}_{X \coprod_Z Y}([e]_{E})\bigr)
=
<^{\mu}_{X \coprod_Z Y}
\bigl(\kappa^{V}_{X \coprod_Z Y}([v]_{V})\bigr).
\]
Consider the source case.
Every such $[e]_E$ is a singleton class, since $R_E$ identifies only
component edges: by the definition of $S_E$ the only class with more than
one element is
$\{\iota^E_X(f_E(\top_Z)), \iota^E_Y(g_E(\top_Z))\}$, and both of its
elements are component edges, since $\top_Z$ is one and $f$ and $g$
preserve labels.
Without loss of generality, write
\[
[e]_E=\{\iota_X^E(e')\}
\]
for some $e'\in (E_X)_\Sigma\cup(E_X)_{\boxsortsubscript}$; in particular
$e' \neq \top_X$, so $<^\mu_X$ is defined at $\kappa^E_X(e')$.
By the definition of the quotient source map,
\[
s_{X\coprod_ZY}([e]_E)
=
([-]_V)^*
\bigl((\iota_X^V)^*(s_X(e'))\bigr).
\]
Hence, if $[v]_V$ occurs in $s_{X\coprod_ZY}([e]_E)$, there exists $v'$
occurring in $s_X(e')$ such that
\[
[\iota_X^V(v')]_V=[v]_V.
\]
Since $X$ satisfies incidence,
\[
<^\mu_{X}\bigl(\kappa^E_{X}(e')\bigr)
=
<^\mu_{X}\bigl(\kappa^V_{X}(v')\bigr),
\]
and applying $\iota^E_X$ to both sides,
\[
<^\mu_{X\coprod Y}
\bigl(\kappa^E_{X\coprod Y}(\iota_X^E(e'))\bigr)
=
<^\mu_{X\coprod Y}
\bigl(\kappa^V_{X\coprod Y}(\iota_X^V(v'))\bigr).
\]
Therefore, by the well-definedness of $<^\mu_{X \coprod Y}/\sim$, whose
value at $[v]_V$ may be computed from the representative
$\iota^V_X(v')$,
\begin{align*}
<^\mu_{X\coprod_ZY}
\bigl(\kappa^V_{X\coprod_ZY}([v]_V)\bigr)
&=
\left[
<^\mu_{X\coprod Y}
\bigl(\kappa^V_{X\coprod Y}(\iota_X^V(v'))\bigr)
\right]_E\\
&=
\left[
<^\mu_{X\coprod Y}
\bigl(\kappa^E_{X\coprod Y}(\iota_X^E(e'))\bigr)
\right]_E\\
&=
<^\mu_{X\coprod_ZY}
\bigl(\kappa^E_{X\coprod_ZY}([e]_E)\bigr).
\end{align*}
The case in which the representative lies in the $Y$-summand is symmetric,
and the target case is analogous.
\item We check non-degeneracy.
Let $[b]_{E} \in (E_{X \coprod_{Z} Y})_{\boxsortsubscript}$.
As noted above, every box edge class is a singleton, so
$[b]_{E} = \{\iota^{E}_{X}(b')\}$ for some $b' \in (E_{X})_{\boxsortsubscript}$,
the case of the $Y$-summand being symmetric.
Since $X$ is a rooted e-hypergraph, $b'$ has at least two children, say
$c_{1} \neq c_{2}$ with $<^{\mu}_{X}(\kappa^{E}_{X}(c_{i})) = b'$; by
alternation in $X$, both are component edges.
Applying $\iota^{E}_{X}$,
\[
<^{\mu}_{X \coprod Y}\bigl(\kappa^{E}_{X \coprod Y}(\iota^{E}_{X}(c_{i}))\bigr)
= \iota^{E}_{X}(b')~,
\]
so $<^{\mu}_{X \coprod_{Z} Y}(\kappa^{E}_{X \coprod_{Z} Y}([\iota^{E}_{X}(c_{i})]_{E})) = [b]_{E}$
for $i = 1, 2$, and $[\iota^{E}_{X}(c_{1})]_{E}$, $[\iota^{E}_{X}(c_{2})]_{E}$
are children of $[b]_{E}$.

It remains to check that they are distinct.
Suppose $[\iota^{E}_{X}(c_{1})]_{E} = [\iota^{E}_{X}(c_{2})]_{E}$.
Since $c_{1} \neq c_{2}$, this class has more than one element, so by the
definition of $S_{E}$ it is
$\{\iota^{E}_{X}(f_{E}(\top_{Z})), \iota^{E}_{Y}(g_{E}(\top_{Z}))\}$, whose two
elements lie in different summands, contradicting
$\iota^{E}_{X}(c_{1}), \iota^{E}_{X}(c_{2}) \in \iota^{E}_{X}(E_{X})$.
Hence $[b]_{E}$ has at least two children.
\end{enumerate}
\end{proof}

\begin{lemma}\label{lemma:pushout_injections_are_homomorphisms}
Define maps
\[
\widetilde{\iota_X}\colon X\longrightarrow X\coprod_ZY,
\qquad
\widetilde{\iota_Y}\colon Y\longrightarrow X\coprod_ZY
\]
by
\begin{align*}
(\widetilde{\iota_X})_V(v)
&=[\iota_X^V(v)]_V,
&
(\widetilde{\iota_X})_E(e)
&=[\iota_X^E(e)]_E,\\
(\widetilde{\iota_Y})_V(v)
&=[\iota_Y^V(v)]_V,
&
(\widetilde{\iota_Y})_E(e)
&=[\iota_Y^E(e)]_E.
\end{align*}
Then $\widetilde{\iota_X}$ is an anchored homomorphism and
$\widetilde{\iota_Y}$ is a strict homomorphism. Moreover, the square
commutes:
\[
\widetilde{\iota_X}\circ f
=
\widetilde{\iota_Y}\circ g.
\]
If $g$ is strict, then $\widetilde{\iota_X}$ is also strict.
\end{lemma}
\begin{proof}
We need to check conditions (1)--(3) of~\cref{def:rooted_homomorphism}.
\begin{enumerate}
\item It must be the case that
\[
((\widetilde{\iota_X})_V)^{*}(s_X(e))
=
s_{X\coprod_ZY}((\widetilde{\iota_X})_E(e))
\]
and
\[
((\widetilde{\iota_X})_V)^{*}(t_X(e))
=
t_{X\coprod_ZY}((\widetilde{\iota_X})_E(e)).
\]
By the definition of the quotient source map~(\ref{def:quotient_source})
and by the definition of $(\widetilde{\iota_X})_V$, we have
\begin{align*}
s_{X\coprod_ZY}((\widetilde{\iota_X})_E(e))
&=
s_{X\coprod_ZY}([\iota_X^E(e)]_E)\\
&=
([-]_V)^*
\bigl((\iota_X^V)^*(s_X(e))\bigr)\\
&=
((\widetilde{\iota_X})_V)^*(s_X(e)).
\end{align*}
Therefore,
\[
((\widetilde{\iota_X})_V)^*(s_X(e))
=
s_{X\coprod_ZY}((\widetilde{\iota_X})_E(e)).
\]
The target case is analogous, and the case of
$\widetilde{\iota_Y}$ is symmetric.

\item It must be the case that
\[
l_X(e)
=
l_{X\coprod_ZY}((\widetilde{\iota_X})_E(e)).
\]
By the definition of the quotient label function~(\ref{def:quotient_label}),
\begin{align*}
l_{X\coprod_ZY}((\widetilde{\iota_X})_E(e))
&=
l_{X\coprod_ZY}([\iota_X^E(e)]_E)\\
&=
l_{X\coprod Y}(\iota_X^E(e))\\
&=
l_X(e).
\end{align*}
The case of $\widetilde{\iota_Y}$ is symmetric.

\item It must be the case that if $e<_X x$, then
\[
(\widetilde{\iota_X})_E(e)
<
_{X\coprod_ZY}
\widetilde{\iota_X}(x).
\]
In fact, we have the stronger property that if $e<^\mu_X x$, then
\[
(\widetilde{\iota_X})_E(e)
<
^\mu_{X\coprod_ZY}
\widetilde{\iota_X}(x).
\]

First, suppose that $x=\kappa_X^V(v)$ for some $v\in V_X$.
Saying that $e<^\mu_X v$ is equivalent to saying that
\[
<^\mu_X(\kappa_X^V(v))=e.
\]
By the definition of $<^\mu_{X\coprod_ZY}$~(\ref{def:quotient_parent}),
we have
\begin{align*}
<^\mu_{X\coprod_ZY}
\bigl(
\kappa^V_{X\coprod_ZY}([\iota_X^V(v)]_V)
\bigr)
&=
\left[
<^\mu_{X\coprod Y}
\bigl(
\kappa^V_{X\coprod Y}(\iota_X^V(v))
\bigr)
\right]_E\\
&=
[\iota_X^E(e)]_E\\
&=
(\widetilde{\iota_X})_E(e).
\end{align*}

The edge case is analogous. If $e<^\mu_X e'$ for some $e'\in E_X$,
then
\begin{align*}
<^\mu_{X\coprod_ZY}
\bigl(
\kappa^E_{X\coprod_ZY}([\iota_X^E(e')]_E)
\bigr)
&=
\left[
<^\mu_{X\coprod Y}
\bigl(
\kappa^E_{X\coprod Y}(\iota_X^E(e'))
\bigr)
\right]_E\\
&=
[\iota_X^E(e)]_E\\
&=
(\widetilde{\iota_X})_E(e).
\end{align*}
Thus $\widetilde{\iota_X}$ preserves every immediate-parent relation.
Consequently, it also preserves their transitive closure, and is therefore
an anchored homomorphism.

The case of $\widetilde{\iota_Y}$ is symmetric. Moreover, we have
\[
(\widetilde{\iota_Y})_E(\top_Y)
=
[\iota_Y^E(\top_Y)]_E
=
\top_{X\coprod_ZY},
\]
so $\widetilde{\iota_Y}$ is strict.

When moreover $g$ is strict, using also the strictness of $f$, we have
\begin{align*}
(\widetilde{\iota_X})_E(\top_X)
&=
[\iota_X^E(\top_X)]_E\\
&=
[\iota_X^E(f_E(\top_Z))]_E\\
&=
[\iota_Y^E(g_E(\top_Z))]_E\\
&=
[\iota_Y^E(\top_Y)]_E\\
&=
\top_{X\coprod_ZY},
\end{align*}
so $\widetilde{\iota_X}$ is also strict.
\end{enumerate}

Finally, the square commutes by the definitions of $R_V$ and $R_E$.
Indeed, for every $v\in V_Z$,
\[
(\widetilde{\iota_X})_V(f_V(v))
=
[\iota_X^V(f_V(v))]_V
=
[\iota_Y^V(g_V(v))]_V
=
(\widetilde{\iota_Y})_V(g_V(v)),
\]
and, since $E_Z=\{\top_Z\}$,
\[
(\widetilde{\iota_X})_E(f_E(\top_Z))
=
[\iota_X^E(f_E(\top_Z))]_E
=
[\iota_Y^E(g_E(\top_Z))]_E
=
(\widetilde{\iota_Y})_E(g_E(\top_Z)).
\]
Therefore,
\[
\widetilde{\iota_X}\circ f
=
\widetilde{\iota_Y}\circ g.
\]
\end{proof}

\begin{lemma}[Universal property]\label{lemma:pushout_universal}
Let $Q$ be a rooted e-hypergraph and let $j_{X} : X \to Q$, $j_{Y} : Y \to Q$ be
lax homomorphisms with $f \fatsemi j_{X} = g \fatsemi j_{Y}$.
Then there is a unique lax homomorphism $u : X \coprod_{Z} Y \to Q$ such that
$\widetilde{\iota_{X}} \fatsemi u = j_{X}$ and $\widetilde{\iota_{Y}} \fatsemi u = j_{Y}$.
If $j_{X}$ and $j_{Y}$ are moreover anchored, then so is $u$; and if $j_{Y}$ is additionally strict, then so is $u$.
\end{lemma}
\begin{proof}
Define $u$ on vertices and edges by
\begin{align*}
  u_{V}([v]_{V}) &= \overline{u_{V}}(v) = [j_{X}^{V},j_{Y}^{V}](v)\\
  u_{E}([e]_{E}) &= \overline{u_{E}}(e) = [j_{X}^{E}, j_{Y}^{E}](e)\\
\end{align*}
First, check that $u$ is well-defined.
For vertices, we need to check that if $[v_1] = [v_2]$, then $\overline{u_{V}}(v_1) = \overline{u_{V}}(v_2)$.
By definition of the equivalence relation $R_{V}$, this means that there exists a sequence $v_1 = w_1 \mathcal{S} \ldots \mathcal{S} w_n = v_2$ such that for each $w_{i} \mathcal{S} w_{i + 1}$, either $w_{i} = w_{i+1}$ or $w_{i} S_{V} w_{i+1}$, or $w_{i+1} S_{V} w_{i}$.
We will do induction on the length of this sequence.
When the length is $2$, the interesting case is when $v_1 S_{V} v_2$.
Without loss of generality, suppose that $v_1 = \iota_{X}^{V}(f_{V}(z))$ and $v_2 = \iota_{Y}^{V}(g_{V}(z))$ for some $z \in V_{Z}$.
Then $\overline{u_{V}}(v_1) = j_{X}^{V}(f_{V}(z)) = j_{Y}^{V}(g_{V}(z)) = \overline{u_{V}}(v_2)$, because $f \fatsemi j_{X} = g \fatsemi j_{Y}$.
The rest follows by symmetry and transitivity.
The case of edges is analogous.

Now we need to check that $\widetilde{\iota_{X}} \fatsemi u = j_{X}$ and $\widetilde{\iota_{Y}} \fatsemi u = j_{Y}$.
For vertices, we have $u(\widetilde{\iota_{X}}(v)) = u([\iota_{X}^{V}(v)]_{V}) = [j_{X}^{V},j_{Y}^{V}](\iota_{X}^{V}(v)) = j_{X}^{V}(v)$.
Similarly for $u(\widetilde{\iota_{Y}}(v))$ and for edges.
Note that we also have
\[
(\iota_{X}^{V} \fatsemi \overline{u_{V}})(v) = \overline{u_{V}}(\iota_{X}^{V}(v)) = u_{V}([\iota_{X}^{V}(v)]_{V}) = u_{V}(\widetilde{\iota_{X}}(v)) = j_{X}^{V}(v)~,
\]
for all $v \in V_{X}$, and hence $\iota_{X}^{V} \fatsemi \overline{u_{V}} = j_{X}^{V}$, and similarly for edges and for $\iota_{Y}$.

Let us check that $u$ is a lax homomorphism.
We will go through the three conditions of~\cref{def:rooted_homomorphism}.
\begin{enumerate}
\item We need to check that
\[
u_V^*(s_{X\coprod_ZY}([e]_E))
=
s_Q(u_E([e]_E))
\]
and similarly for targets.
By the definition of the quotient source map,
\[
s_{X\coprod_ZY}([e]_E)
=
([-]_V)^*(s_{X\coprod Y}(e)).
\]
Therefore, since $[-]_V\fatsemi u_V=\overline u_V$,
\begin{align*}
u_V^*(s_{X\coprod_ZY}([e]_E))
&=
u_V^*\bigl(([-]_V)^*(s_{X\coprod Y}(e))\bigr)\\
&=
\overline u_V^*(s_{X\coprod Y}(e)).
\end{align*}
Without loss of generality, suppose that
$e=\iota_X^E(e')$. Then
\[
s_{X\coprod Y}(e)
=
(\iota_X^V)^*(s_X(e')).
\]
Since $\iota_X^V\fatsemi\overline u_V=j_{X}^{V}$, we obtain
\begin{align*}
\overline u_V^*(s_{X\coprod Y}(e))
&=
\overline u_V^*\bigl((\iota_X^V)^*(s_X(e'))\bigr)\\
&=
(j_{X}^{V})^*(s_X(e'))\\
&=
s_Q(j_{X}^{E}(e')),
\end{align*}
where the last equality follows because $j_X$ is a lax homomorphism.
On the other hand,
\[
u_E([e]_E)
=
\overline u_E(e)
=
\overline u_E(\iota_X^E(e'))
=
j_{X}^{E}(e').
\]
Consequently,
\[
u_V^*(s_{X\coprod_ZY}([e]_E))
=
s_Q(u_E([e]_E)).
\]
The case in which $e = \iota_{Y}^{E}(e')$ is symmetric, and the
target calculation is analogous.
\item We now check that
\[
l_{X\coprod_ZY}([e]_E)=l_Q(u_E([e]_E)).
\]
By the definition of the quotient label map,
\[
l_{X\coprod_ZY}([e]_E)=l_{X\coprod Y}(e).
\]
Without loss of generality, suppose that $e=\iota_X^E(e')$. Then
\[
l_{X\coprod Y}(e)=l_X(e').
\]
Since $j_X$ is a lax homomorphism,
\[
l_X(e')=l_Q(j_X^E(e')).
\]
On the other hand,
\[
u_E([e]_E)
=
\overline u_E(e)
=
\overline u_E(\iota_X^E(e'))
=
j_X^E(e').
\]
Hence
\[
l_{X\coprod_ZY}([e]_E)=l_Q(u_E([e]_E)).
\]
The case in which $e = \iota_{Y}^{E}(e')$ is symmetric.
\item We now check preservation of the hierarchical relation. We need
to show that
\[
[e]_E<_{X\coprod_Z Y}[v]_V
\quad\Longrightarrow\quad
u_E([e]_E)<_Q u_V([v]_V)~,
\]
and that
\[
[e]_E<_{X\coprod_Z Y}[e']_E
\quad\Longrightarrow\quad
u_E([e]_E)<_Q u_E([e']_E)~.
\]
By definition, $<_{X \coprod_{Z} Y}$ is the transitive closure of the immediate predecessor relation $<^{\mu}_{X \coprod_{Z} Y}$.
Consider the case of vertices, and let $[v]_{V} \in V_{X \coprod_{Z} Y}$.
Suppose that $[e]_{E} <_{X \coprod_{Z} Y} [v]_{V}$.
This means that there is a chain of edges in $X \coprod_{Z} Y$ such that $[e]_{E} <^{\mu}_{X \coprod_{Z} Y} [e_1]_{E} <^{\mu}_{X \coprod_{Z} Y} \ldots <^{\mu}_{X \coprod_{Z} Y} [e_n]_{E} <^{\mu}_{X \coprod_{Z} Y} [v]_{V}$.
It suffices to check that every adjacent pair in this chain is preserved by $u$.
We will show that it holds for the pair $[e]_{E} <^{\mu}_{X \coprod_{Z} Y} [v]_{V}$, then the rest follow by analogous argument and transitivity.
We need to check that $u_{E}([e]_{E}) <_{Q} u_{V}([v]_{V})$.
By definition of $<^{\mu}_{X \coprod_{Z} Y}$, we have that $[e]_{E} = <^{\mu}_{X \coprod_{Z} Y}(\kappa_{X \coprod_{Z} Y}^{V}([v]_{V})) = [<_{X \coprod Y}^{\mu}(\kappa_{X \coprod Y}^{V}(v'))]_{E}$ for some $v'$ such that $<_{X \coprod Y}^{\mu}(\kappa_{X \coprod Y}^{V}(v'))$ is defined.
Without loss of generality, suppose that $v' = \iota_{X}^{V}(v'')$, so that $<_{X \coprod Y}^{\mu}(\kappa_{X \coprod Y}^{V}(\iota_{X}^{V}(v''))) = \iota_{X}^{E}(<_{X}^{\mu}(\kappa_{X}^{V}(v'')))$.
Therefore, $[e]_{E} = [<_{X \coprod Y}^{\mu}(\kappa_{X \coprod Y}^{V}(\iota_{X}^{V}(v'')))]_{E} = [\iota_{X}^{E}(<_{X}^{\mu}(\kappa_{X}^{V}(v'')))]_{E}$.
Since $j_{X}$ is a lax homomorphism, we have that $j_{X}^{E}(<_{X}^{\mu}(\kappa_{X}^{V}(v''))) <_{Q} j_{X}^{V}(v'')$.
By definition of $u$, and because the value of $u$ does not depend on the choice of a representative of $[e]_{E}$, we may pick $e = \iota_{X}^{E}(<_{X}^{\mu}(\kappa_{X}^{V}(v'')))$ and hence 
\begin{align*}
u_{E}([e]_{E}) &= \overline{u_{E}}(e)\\
&= \overline{u_{E}}(\iota_{X}^{E}(<_{X}^{\mu}(\kappa_{X}^{V}(v''))))\\
&= j_{X}^{E}(<_{X}^{\mu}(\kappa_{X}^{V}(v'')))\\
& <_{Q} j_{X}^{V}(v'')\\
&= \overline{u_{V}}(\iota_{X}^{V}(v'')) = u_{V}([v]_{V})~.
\end{align*}
% The rest follows by the analogous argument for $[e_{i}]_{E} <^{\mu}_{X \coprod_{Z} Y} [e_{i+1}]_{E}$ and transitivity.
If $j_{X}$ and $j_{Y}$ are anchored, then $j_{X}^{E}(<_{X}^{\mu}(\kappa_{X}^{V}(v''))) <^{\mu}_{Q} j_{X}^{V}(v'')$ and $j_{Y}^{E}(<_{Y}^{\mu}(\kappa_{Y}^{V}(v''))) <^{\mu}_{Q} j_{Y}^{V}(v'')$, therefore $u_{E}([e]_{E}) <^{\mu}_{Q} u_{V}([v]_{V})$ (by the chain above) and $u$ is anchored.
Moreover, if $j_{Y}$ is strict, then $u_{E}(\top_{X \coprod_{Z} Y}) = u_{E}([\iota_{Y}^{E}(\top_{Y})]_{E}) = j_{Y}^{E}(\top_{Y}) = \top_{Q}$.
\end{enumerate}
It remains to check that $u$ is unique.
Let
\[
u':X\coprod_Z Y\to Q
\]
be another homomorphism such that
\[
\widetilde{\iota_X}\fatsemi u'=j_X,
\qquad
\widetilde{\iota_Y}\fatsemi u'=j_Y.
\]
Let $[v]_V\in V_{X\coprod_ZY}$. 
This class has a representative in either the $X$- or the $Y$-summand. Without loss of generality, suppose
that
\[
[v]_V=(\widetilde{\iota_X})_V(v')
\]
for some $v'\in V_X$. Then
\begin{align*}
u'_V([v]_V)
&=
u'_V((\widetilde{\iota_X})_V(v'))\\
&=
j_X^V(v')\\
&=
u_V((\widetilde{\iota_X})_V(v'))\\
&=
u_V([v]_V).
\end{align*}
The case in which $[v]_V$ is represented in the $Y$-summand is
analogous. Hence $u'_V=u_V$. The same argument for edge classes gives
$u'_E=u_E$. Therefore $u'=u$.
\end{proof}

\begin{example}
  Consider the diagram below
  \[
  \adjustbox{max width=0.65\textwidth}{
    \tikzfig{./figures/monoidal_egraphs/pushout_example}
  }~,
  \]
  where node and edge labels together with dashed arrows suggest how maps $f$ and $g$ are defined.
  The rooted e-hypergraph $X \coprod_{Z} Y$ along with $\widetilde{\iota_{X}}$ and $\widetilde{\iota_{Y}}$ is then the pushout for the square.
  The underlying principle is simple: since $f$ is strict and $g$ is anchored, after the images of $x_1,x_2$ and $y_1,y_2$ were identified to make the square commute, the image of $\top_{X}$ is identified with the image of $\cmpsort$-edge $Y_1$.
  Intuitively, it corresponds to gluing $X$ and $Y$ along the common interface $Z$ at the level specified by $g$---requiring $f$ to be strict makes such gluing unambiguous, unlike in~\cref{example:pushout_failure}.
\end{example}

With the pushout construction in hand, we can proceed to defining the category of rooted e-hypergraphs with an appropriate notion of interface.

\begin{definition}
Given pointed discrete ordered $n$ and $m$, we write $n \sqcup m$ for the
pushout of $n \leftarrow \mathfrak{T} \rightarrow m$ in
$\catname{EHyp}_{\mathrm{s}}(\Sigma)$, which identifies the two roots and
concatenates the vertex orders, and $n \setminus m$ for the pointed discrete
e-hypergraph obtained by removing the vertices of $m$ when $m$ is a
sub-e-hypergraph of $n$.
% An \emph{interface leg} is a lax homomorphism $\phi : n \to \mathcal{G}$ from
% a pointed discrete $n$ such that $\phi(\top_{n}) = \top_{\mathcal{G}}$.
\end{definition}

These pointed discrete ordered e-hypergraphs will serve as interfaces.
We can equivalently describe pointed discrete ordered e-hypergraphs by using a functor $D : \mathbb{F} \to \catname{EHyp}_{\mathrm{s}}(\Sigma)$, similarly to~\cref{def:hypI}.

\begin{definition}\label{def:discrete_functor}
Let $\mathbb{F}$ be the skeleton of the category of finite sets, with objects
$[n] = \{1, \ldots, n\}$.
Define $D : \mathbb{F} \to \catname{EHyp}_{\mathrm{s}}(\Sigma)$ by letting
$D[n]$ be the pointed discrete e-hypergraph with vertices
$v_{1}, \ldots, v_{n}$ ordered by $v_{1} < \cdots < v_{n}$, whose only edge is
$\top$ with $l(\top) = \cmpsort$ and $\top <^{\mu} v_{i}$ for every $i$, and by
letting $D(\varphi) : D[n] \to D[m]$, for $\varphi : [n] \to [m]$, be given by
$v_{i} \mapsto v_{\varphi(i)}$ and $\top \mapsto \top$.
Given pointed discrete ordered $n$ and $m$, we write $n \sqcup m$ for the
pointed discrete ordered e-hypergraph obtained by taking the disjoint union of
their vertices, identifying their roots, and ordering the result so that the
vertices of $n$ precede those of $m$; and $n \setminus m$ for the pointed
discrete e-hypergraph obtained by removing the vertices of $m$ from $n$, when
$m$ is a sub-e-hypergraph of $n$.
\end{definition}

\begin{remark}\label{remark:discrete_functor_strict}
Each $D(\varphi)$ is indeed a strict homomorphism, with no assumption on
$\varphi$: condition (1) of~\cref{def:rooted_homomorphism} is vacuous because
$\cmpsort$-edges are $0$-ary, condition (2) holds as both roots are
$\cmpsort$-labelled, and condition (3$'$) holds because every vertex of a
pointed discrete e-hypergraph is an immediate child of its root.
Moreover $n \sqcup m$ is the pushout of
$n \xleftarrow{} \mathfrak{T} \xrightarrow{} m$ in
$\catname{EHyp}_{\mathrm{s}}(\Sigma)$: the span satisfies the hypotheses
of~[the pushout proposition], since $\mathfrak{T}$ is pointed discrete and both
legs are strict, and the construction there identifies the two roots and leaves
everything else untouched, as $\mathfrak{T}$ has no vertices.
The order is not part of the pushout data; we fix it by the convention above,
so that the vertices of $n$ come before those of $m$.
Consequently $D[n] \sqcup D[m] = D[n+m]$, and $D$ is a strict symmetric
monoidal functor from $(\mathbb{F}, +, [0])$ to
$(\catname{EHyp}_{\mathrm{s}}(\Sigma), \sqcup, \mathfrak{T})$.
\end{remark}
% \begin{remark}
%   Just as in~\cref{def:hypI} we can consider a functor $D : \mathbb{F} \to \catname{EHyp}_{\leq}(\Sigma)$ that maps $[n]$ to a pointed discrete ordered e-hypergraph with $n$ vertices and $\top$ as the only edge.
%   $[n] + [m]$ is then mapped to the fibered coproduct, i.e., a pushout along a span $D([n]) \xleftarrow{} \top \xrightarrow{} D([m])$.
%   This construction says that in the fibered coproduct vertices of $D([n])$ come before those of $D([m])$.
% \end{remark}

\begin{remark}\label{remark:compare_interfaces}
From a pointed discrete domain, condition (3) of
\cref{def:rooted_homomorphism} requires only that vertex images be proper
descendants of the root, which excludes nothing: provided that the $\top$ element of the pointed discrete e-hypergraph is mapped to the $\top$ element of another e-hypergraph, vertices may then be freely mapped without any restrictions. 
This matches the effectively unconstrained
interface legs of~\cite[Def.~IV.5]{tiurin2025equivalencehypergraphsdporewriting}.
An interface leg is strict precisely when all images are immediate children of
$\top_{\mathcal{G}}$; this replaces the requirement of op.\ cit.\ that the
images of external interfaces consist of top-level vertices.
\end{remark}

% \begin{remark}
% Equivalently, for a pointed discrete $n$ one could consider a functor $D : \mathbb{F} \to \catname{EHyp}(\Sigma)$ as in~\cref{def:hypI} that would map a coproduct in $\mathbb{F}$ to a fibered coproduct in $\catname{EHyp}(\Sigma)$ and map a set $[n]$ to a rooted e-hypergraph with $n$ vertices such that $\top$ is the immediate parent of all vertices.
% \end{remark}
%% ---------------------------------------------------------------------------
%% Extended cospans
%% ---------------------------------------------------------------------------

\begin{definition}
  An \emph{extended cospan} is a diagram in $\catname{EHyp}_{\leq}(\Sigma)$ of the form
  \[
  n \xrightarrow{f_{\mathrm{ext}}} n' \xrightarrow{f_{\mathrm{int}}} \mathcal{G} \xleftarrow{g_{\mathrm{int}}} m' \xleftarrow{g_{\mathrm{ext}}} m~,
  \]
  such that 
  \begin{itemize}
  \item $n,n',m,m'$ are pointed discrete ordered e-hypergraphs;
  \item $f_{\mathrm{ext}}$ and $g_{\mathrm{ext}}$ are mono;
  \item $f_{\mathrm{ext}} \fatsemi f_{\mathrm{int}}$ and $g_{\mathrm{ext}} \fatsemi g_{\mathrm{int}}$ are strict;
  % \item The images of $f_{\mathrm{int},V} \restriction_{n' \setminus f_{\mathrm{ext}}(n)}$ and $g_{\mathrm{int},V} \restriction_{m' \setminus g_{\mathrm{ext}}(m)}$ consist of exclusively of vertices whose immediate predecessors are not $\top_{\mathcal{G}}$.
  \item
  \[
  f_{\mathrm{int},V}
  \bigl(V_{n'} \setminus f_{\mathrm{ext},V}(V_n)\bigr)
  \subseteq
  \left\{
  v \in V_{\mathcal{G}}
  \;\middle|\;
  <^\mu_{\mathcal{G}}(\kappa^V_{\mathcal{G}}(v))
  \neq \top_{\mathcal{G}}
  \right\},
  \]
and similarly,
\[
g_{\mathrm{int},V}
\bigl(V_{m'} \setminus g_{\mathrm{ext},V}(V_m)\bigr)
\subseteq
\left\{
v \in V_{\mathcal{G}}
\;\middle|\;
<^\mu_{\mathcal{G}}(\kappa^V_{\mathcal{G}}(v))
\neq \top_{\mathcal{G}}
\right\}.
\]
  \end{itemize}
  We refer to the vertices in $V_{n'} \setminus f_{\mathrm{ext},V}(V_n)$ and $V_{m'} \setminus g_{\mathrm{ext},V}(V_m)$, together with the corresponding restrictions of $f_{\mathrm{int},V}$ and $g_{\mathrm{int},V}$, as the \emph{strictly internal} interfaces.
\end{definition}

The last condition above ensures that the image of strictly internal interfaces consists exclusively of vertices nested inside $\cmpsort$-edges.
Note that the external interfaces are always strict.

% \begin{definition}[Isomorphism of extended cospans.]\label{def:isomorphic_ecospans}
% Consider two extended cospans below
% % https://q.uiver.app/#q=WzAsOCxbMCwxLCJuIl0sWzEsMCwibiciXSxbMSwyLCJuJyciXSxbMywwLCJcXG1hdGhjYWx7R30iXSxbMywyLCJcXG1hdGhjYWx7Ryd9Il0sWzUsMCwibSciXSxbNSwyLCJtJyciXSxbNiwxLCJtIl0sWzAsMSwiIiwwLHsiY3VydmUiOi0yfV0sWzAsMiwiIiwyLHsiY3VydmUiOjJ9XSxbMSwzXSxbMiw0XSxbNSwzXSxbNiw0XSxbNyw1LCIiLDIseyJjdXJ2ZSI6Mn1dLFs3LDYsIiIsMCx7ImN1cnZlIjotMn1dLFsxLDIsIlxcYWxwaGEiXSxbMyw0LCJcXGJldGEiXSxbNSw2LCJcXGdhbW1hIl1d
% \[\begin{tikzcd}
% 	& {n'} && {\mathcal{G}} && {m'} & \\
% 	n &&&&&& m \\
% 	& {n''} && {\mathcal{G'}} && {m''}
% 	\arrow[from=1-2, to=1-4]
% 	\arrow["\alpha", from=1-2, to=3-2]
% 	\arrow["\beta", from=1-4, to=3-4]
% 	\arrow[from=1-6, to=1-4]
% 	\arrow["\gamma", from=1-6, to=3-6]
% 	\arrow[curve={height=-12pt}, from=2-1, to=1-2]
% 	\arrow[curve={height=12pt}, from=2-1, to=3-2]
% 	\arrow[curve={height=12pt}, from=2-7, to=1-6]
% 	\arrow[curve={height=-12pt}, from=2-7, to=3-6]
% 	\arrow[from=3-2, to=3-4]
% 	\arrow[from=3-6, to=3-4]
% \end{tikzcd}\]
% and consider a relation
% \[
% R = \{ (x,y) \mid z <^{\mu} f_{\mathrm{int}}(x) \land z <^{\mu} f_{\mathrm{int}}(y), \quad x,y \in n' \}
% \]
% and let $S$ be its reflexive closure. The latter partitions $n'$ into non-empty subsets---each indexed by a $\cmpsort$-edge of $\mathcal{G}$.
% We get analogous partitions for $m'$, $n''$, and $m''$.
% Two extended cospans are isomorphic if there exist isomorphisms $\alpha : n' \to n''$, $\beta : \mathcal{G} \to \mathcal{G'}$, and $\gamma : m' \to m''$ that make the above diagram commute, and such that $\alpha$ and $\gamma$ preserve the order within each block of the partition.
% \end{definition}
\begin{definition}[Isomorphism of extended cospans.]\label{def:isomorphic_ecospans}
Consider two extended cospans with the same feet $n$ and $m$,
\[
n \xrightarrow{f_{\mathrm{ext}}} n' \xrightarrow{f_{\mathrm{int}}} \mathcal{G}
  \xleftarrow{g_{\mathrm{int}}} m' \xleftarrow{g_{\mathrm{ext}}} m~,
\qquad
n \xrightarrow{h_{\mathrm{ext}}} n'' \xrightarrow{h_{\mathrm{int}}} \mathcal{G}'
  \xleftarrow{k_{\mathrm{int}}} m'' \xleftarrow{k_{\mathrm{ext}}} m~,
\]
as shown in the diagram below.
% https://q.uiver.app/#q=WzAsOCxbMCwxLCJuIl0sWzEsMCwibiciXSxbMSwyLCJuJyciXSxbMywwLCJcXG1hdGhjYWx7R30iXSxbMywyLCJcXG1hdGhjYWx7Ryd9Il0sWzUsMCwibSciXSxbNSwyLCJtJyciXSxbNiwxLCJtIl0sWzAsMSwiIiwwLHsiY3VydmUiOi0yfV0sWzAsMiwiIiwyLHsiY3VydmUiOjJ9XSxbMSwzXSxbMiw0XSxbNSwzXSxbNiw0XSxbNyw1LCIiLDIseyJjdXJ2ZSI6Mn1dLFs3LDYsIiIsMCx7ImN1cnZlIjotMn1dLFsxLDIsIlxcYWxwaGEiXSxbMyw0LCJcXGJldGEiXSxbNSw2LCJcXGdhbW1hIl1d
\[\begin{tikzcd}
	& {n'} && {\mathcal{G}} && {m'} & \\
	n &&&&&& m \\
	& {n''} && {\mathcal{G'}} && {m''}
	\arrow["{f_{\mathrm{int}}}", from=1-2, to=1-4]
	\arrow["\alpha", from=1-2, to=3-2]
	\arrow["\beta", from=1-4, to=3-4]
	\arrow["{g_{\mathrm{int}}}"', from=1-6, to=1-4]
	\arrow["\gamma", from=1-6, to=3-6]
	\arrow["{f_{\mathrm{ext}}}", curve={height=-12pt}, from=2-1, to=1-2]
	\arrow["{h_{\mathrm{ext}}}"', curve={height=12pt}, from=2-1, to=3-2]
	\arrow["{g_{\mathrm{ext}}}", curve={height=12pt}, from=2-7, to=1-6]
	\arrow["{k_{\mathrm{ext}}}"', curve={height=-12pt}, from=2-7, to=3-6]
	\arrow["{h_{\mathrm{int}}}"', from=3-2, to=3-4]
	\arrow["{k_{\mathrm{int}}}", from=3-6, to=3-4]
\end{tikzcd}\]
By alternation, the immediate parent of a vertex is a $\cmpsort$-edge, so we
may consider the map
\[
\pi_{f} : n' \to (E_{\mathcal{G}})_{\cmpsortsubscript}~,
\qquad
x \longmapsto {<^{\mu}_{\mathcal{G}}}\bigl(f_{\mathrm{int}}(x)\bigr)~.
\]
Its non-empty preimages
$\pi_{f}^{-1}(e) = \{ x \in n' \mid \pi_{f}(x) = e \}$ partition $n'$ into
non-empty subsets, which we call \emph{blocks}---each indexed by a
$\cmpsort$-edge of $\mathcal{G}$.
The block indexed by $\top_{\mathcal{G}}$ is exactly $f_{\mathrm{ext}}(n)$, and
the remaining blocks partition the strictly internal interface
$n' \setminus f_{\mathrm{ext}}(n)$.
We define $\pi_{g}$, $\pi_{h}$ and $\pi_{k}$ analogously, and obtain in the
same way partitions of $m'$, $n''$ and $m''$.

Two extended cospans are isomorphic if there exist isomorphisms
$\alpha : n' \to n''$ and $\gamma : m' \to m''$ of pointed discrete
e-hypergraphs and an isomorphism $\beta : \mathcal{G} \to \mathcal{G}'$ in
$\catname{EHyp}_{\mathrm{s}}(\Sigma)$ that make the above diagram commute, and
such that $\alpha$ and $\gamma$ preserve the order within each block other than
the one indexed by $\top_{\mathcal{G}}$, respectively $\top_{\mathcal{G}'}$.
\end{definition}


\begin{figure}
\[
\adjustbox{scale=0.55}{
  \tikzfig{./figures/monoidal_egraphs/extended_cospan_example}
}
\]
\caption{Extended cospan example.}
\label{fig:extended_cospan_example}
\end{figure}

\begin{example}\label{example:non_isomorphic_cospans}
Consider the extended cospan in~\cref{fig:extended_cospan_example}.
The labels again suggest how the interface maps $f_{\mathrm{ext}}, f_{\mathrm{int}}$ (respectively, $g_{\mathrm{ext}}, g_{\mathrm{int}}$) are defined.
The partitions of the interfaces induced by the $\cmpsort$-edges $E_1$ and $E_2$ are given with colours.
The notion of the isomorphism from~\cref{def:isomorphic_ecospans} suggests that the position of the image of $n$ within $n'$ relative to the strict internal interface does not matter.
For example, vertices 
\begin{tikzpicture}
	\begin{pgfonlayer}{nodelayer}
		\node [style=red node, label={above:$w_{3}$}] (0) at (-0.5, 3.5) {};
		\node [style=red node, label={above:$w_4$}] (1) at (0, 3.5) {};
	\end{pgfonlayer}
\end{tikzpicture}
can be swapped with vertices
\begin{tikzpicture}
	\begin{pgfonlayer}{nodelayer}
		\node [style=blue node, label={above:$w_{5}$}] (0) at (-0.5, 3.5) {};
		\node [style=blue node, label={above:$w_6$}] (1) at (0, 3.5) {};
	\end{pgfonlayer}
\end{tikzpicture}
in the input interface yielding an isomorphic cospan.
However, we may not swap
\begin{tikzpicture}
	\begin{pgfonlayer}{nodelayer}
		\node [style=red node, label={above:$w_{3}$}] (0) at (-0.5, 3.5) {};
		\node [style=red node, label={above:$w_4$}] (1) at (0, 3.5) {};
	\end{pgfonlayer}
\end{tikzpicture}
with each other---this would yield a non-isomorphic cospan.
\end{example}

\begin{example}
  A pair of non-isomorphic extended cospans is given below
  \[
  \adjustbox{scale=0.55}{
  \tikzfig{./figures/monoidal_egraphs/sym_cospan}
  } \begin{array}{c}
    \xrightarrow{\gamma}\\[8ex]
    \xrightarrow{\beta}\\[8ex]
    \xrightarrow{\alpha}
  \end{array} 
  \quad
\adjustbox{scale=0.55}{
  \tikzfig{./figures/monoidal_egraphs/id_x_id_cospans}
}~.
  \]
  We can try to build an isomorphism between them to see where it fails.
  Let the components of the potential candidate be isomorphisms $\alpha, \beta, \gamma$.
  Note that $v_2$ and $v_3$ of the left cospan's carrier must be mapped by $\beta$ to $v_2$ and $v_3$ of the right cospan's carrier, respectively, so that the triangles $(f_{\mathrm{int}}^{X} \fatsemi \beta)(v_2) = (\alpha \fatsemi f_{\mathrm{int}}^{Y})(w_2)$ and $(f_{\mathrm{int}}^{X} \fatsemi \beta)(v_3) = (\alpha \fatsemi f_{\mathrm{int}}^{Y})(w_3)$.
  On the other hand, from the output interfaces, we have that $(g_{\mathrm{int}}^{X} \fatsemi \beta)(u_2) = \beta(v_3) = g_{\mathrm{int}}^{Y}(\gamma(u_2)) = v_2$ and thus we have a contradiction.
  Intuitively, the left component of the $\boxsort$-edge in the left cospan represents a symmetry morphism, whereas the left component of the $\boxsort$-edge in the right cospan represents a tensor product of two identity morphisms; clearly these should not be identified.

  Below is another pair of non-isomorphic extended cospans
  \[
  \adjustbox{scale=0.55}{
    \tikzfig{./figures/monoidal_egraphs/non_isomorphic_fail_2_left}
  }\qquad
  \adjustbox{scale=0.55}{
    \tikzfig{./figures/monoidal_egraphs/non_isomorphic_fail_2_right}
  }~.
  \]
  These cospans correspond to $\catname{SLat}$ morphism terms $\sym \fatsemi (\sym \join (\sym \fatsemi (f \otimes g)))$ and $(\id \otimes \id) \join (f \otimes g)$, respectively.
  These two terms are equal as morphism in an $\catname{SLat}$-category, however they are not equal as cospans.
  We will equate such cospans at a later stage with rewriting---the equality comes from the interaction with $\catname{SLat}$-structure, so it requires a special treatment.
\end{example}

% \begin{definition}[Well typed extended cospan]
% Consider an extended cospan
% \[
% n \xrightarrow{f_{\mathrm{ext}}} n' \xrightarrow{f_{\mathrm{int}}} \mathcal{G} \xleftarrow{g_{\mathrm{int}}} m' \xleftarrow{g_{\mathrm{ext}}} m~.
% \]
% Take a $\boxsort$ edge $e$ from $\mathcal{G}$ and let $e_1, \ldots, e_k$ be its immediate children.
% Note that each $e_i$ is a $\cmpsort$-edge.
% We call a $\cmpsort$-edge $e_i$ \emph{well-typed} if $|\{ x \in V_{n'} \mid e_i <^{\mu} f_{\mathrm{int},V}(x) \}| = |s(e)|$ and $|\{ x \in V_{m'} \mid e_i <^{\mu} g_{\mathrm{int}}(x)\}| = |t(e)|$.
% We call $e$ well-typed if all of its immediate children edges are well-typed.
% Finally, we call an extended cospan well-typed if all of its $\boxsort$-edges are well-typed.
% \end{definition}
\begin{definition}[Well typed extended cospan]
Consider an extended cospan
\[
n \xrightarrow{f_{\mathrm{ext}}} n' \xrightarrow{f_{\mathrm{int}}} \mathcal{G} \xleftarrow{g_{\mathrm{int}}} m' \xleftarrow{g_{\mathrm{ext}}} m~,
\]
and let $\pi_{f}$ and $\pi_{g}$ be as in~\cref{def:isomorphic_ecospans}.
Take a $\boxsort$-edge $e$ from $\mathcal{G}$ and let $e_1, \ldots, e_k$ be its
immediate children; by alternation each $e_i$ is a $\cmpsort$-edge, and by
non-degeneracy $k \geq 2$.
We call a $\cmpsort$-edge $e_i$ \emph{well-typed} if
\[
|\pi_{f}^{-1}(e_i)| = |s(e)|
\qquad\text{and}\qquad
|\pi_{g}^{-1}(e_i)| = |t(e)|~,
\]
that is, if the block of the internal input interface indexed by $e_i$ has as
many vertices as $e$ has sources, and dually for the outputs.
We call $e$ well-typed if all of its immediate children are well-typed, and an
extended cospan well-typed if all of its $\boxsort$-edges are well-typed.
\end{definition}

Intuitively, well-typedness means that every component edge of a
$\boxsort$-edge has the same number of input and output vertices as the
$\boxsort$-edge itself.
The extended cospan in~\cref{fig:extended_cospan_example} is well-typed.
The $\boxsort$-edge $E_3$ has two sources and two targets; each of its
components $E_1$ and $E_2$ indexes a block of two vertices in the internal
input interface, and likewise a block of two vertices in the internal output
interface.

Well-typedness also ensures that the internal interfaces determine how each
component is wired to the ports of its $\boxsort$-edge.
For a component $e_i$ of a $\boxsort$-edge $e$, the block $\pi_{f}^{-1}(e_i)$
and the sequence $s(e)$ have the same length, so matching them in order yields
a bijection recording which port of $e$ each vertex of the block stands for;
dually for $\pi_{g}^{-1}(e_i)$ and $t(e)$.
In~\cref{fig:extended_cospan_example}, the block $\pi_{f}^{-1}(E_1)$ is
$w_3, w_4$, so $w_3$ stands for the first source $w_1$ of $E_3$ and $w_4$ for
the second source $w_2$; the block $\pi_{f}^{-1}(E_2)$ is $w_5, w_6$, which
stand for $w_1$ and $w_2$ respectively.
This bijection is what makes the order within a strictly internal block part of
the data, whereas the order between blocks is not.
The latter is what allows us to eventually absorb the commutativity of the $\catname{SLat}$-structure.

% \begin{definition}[Category of rooted extended cospans]\label{def:rooted_ehypi}
% The category $\catname{EHypI}_{\bullet}(\Sigma)$ has as objects pointed
% discrete ordered e-hypergraphs, and as morphisms $m \to n$ the isomorphism
% classes of \emph{well typed extended cospans}
% \[
% n \xrightarrow{f_{\mathrm{ext}}} n' \xrightarrow{f_{\mathrm{int}}} \mathcal{G}
%   \xleftarrow{g_{\mathrm{int}}} m' \xleftarrow{g_{\mathrm{ext}}} m~.
% \]

% Composition of two morphisms 
% \begin{align*}
% 	n \xrightarrow{f_{ext}} n' \xrightarrow{f_{int}} &\mathcal{F} \xleftarrow{f'_{int}} m' \xleftarrow{f'_{ext}} m\\
% 	m \xrightarrow{g_{ext}} m'\!' \xrightarrow{g_{int}} &\mathcal{G} \xleftarrow{g'_{int}} k' \xleftarrow{g'_{ext}} k
% \end{align*} is computed in two stages.
% First, $\mathcal{H}$ is computed as the result of the pushout square shown below: 
% % https://q.uiver.app/#q=WzAsMTAsWzAsMCwiZXh0KFgpIl0sWzEsMSwiWCJdLFsyLDIsIkEiXSxbMywwLCJZIl0sWzQsMCwiZXh0KFkpIl0sWzUsMCwiWSciXSxbNiwyLCJCIl0sWzcsMSwiWiJdLFs4LDAsImV4dChaKSJdLFs0LDMsIl57XFx1bGNvcm5lcn1BK197Zl97ZXh0fTtmLGdfe2V4dH07Z31CIl0sWzAsMV0sWzEsMl0sWzMsMiwiZiJdLFs0LDMsImZfe2V4dH0iXSxbNCw1LCJnX3tleHR9IiwyXSxbNSw2LCJnIiwyXSxbNyw2XSxbOCw3XSxbMiw5XSxbNiw5XV0=
% \[
% \trimbox{0cm 0cm 0cm 0.75cm}{
% \adjustbox{scale=0.9,center}
% {\tikzfig{../figures/combinatorial_semantics/pushout_e_cospans}}
% }
% \]
% then,  the result of composition is defined as follows:
% \[
% n \xrightarrow{f_{ext};\iota_1} n' \sqcup (m'' \setminus m) \xrightarrow{h_{1}} \mathcal{H} \xleftarrow{h_2} k' \sqcup (m' \setminus m) \xleftarrow{g'_{ext};\iota_1} k
% \]
% where $h_i$ are defined as follows,  using $|$ to denote the restriction of a function on a discrete e-hypergraph. 
% \[
%     h_1 = [ f_{int};p_1, ~(g_{int};p_2)|_{m'' \setminus m} ]
% \;
%     h_2 = [ g'_{int};p_2, ~(f'_{int};p_1)|_{m' \setminus m} ]
% \]
% The identity of composition is given by the obvious extended cospan of identities. 
% $\catname{EHypI}_{\bullet}(\Sigma)$ inherits a symmetric monoidal structure from the fibered coproduct (and the rooted e-hypergraph $\mathfrak{T}$) structure of $\catname{EHyp({\Sigma}})_{\leq}$.
% \end{definition}

\begin{definition}[Category of rooted extended cospans]\label{def:rooted_ehypi}
The category $\catname{EHypI}_{\bullet}(\Sigma)$ has as objects pointed
discrete ordered e-hypergraphs, and as morphisms $n \to m$ the isomorphism
classes of \emph{well typed extended cospans}
\[
n \xrightarrow{f_{\mathrm{ext}}} n' \xrightarrow{f_{\mathrm{int}}} \mathcal{G}
  \xleftarrow{g_{\mathrm{int}}} m' \xleftarrow{g_{\mathrm{ext}}} m~.
\]
Composition of two morphisms
\begin{align*}
	n \xrightarrow{f_{\mathrm{ext}}} n' \xrightarrow{f_{\mathrm{int}}} &\mathcal{F} \xleftarrow{f'_{\mathrm{int}}} m' \xleftarrow{f'_{\mathrm{ext}}} m\\
	m \xrightarrow{g_{\mathrm{ext}}} m'' \xrightarrow{g_{\mathrm{int}}} &\mathcal{G} \xleftarrow{g'_{\mathrm{int}}} k' \xleftarrow{g'_{\mathrm{ext}}} k
\end{align*}
is computed in two stages.
First, $\mathcal{H}$ is computed as the pushout of the span
\[
\mathcal{F}
  \xleftarrow{\;f'_{\mathrm{ext}} \fatsemi f'_{\mathrm{int}}\;} m
  \xrightarrow{\;g_{\mathrm{ext}} \fatsemi g_{\mathrm{int}}\;} \mathcal{G}~,
\]
with coprojections $p_{1} : \mathcal{F} \to \mathcal{H}$ and
$p_{2} : \mathcal{G} \to \mathcal{H}$, as shown in the square below.
% https://q.uiver.app/#q=WzAsMTAsWzAsMCwiZXh0KFgpIl0sWzEsMSwiWCJdLFsyLDIsIkEiXSxbMywwLCJZIl0sWzQsMCwiZXh0KFkpIl0sWzUsMCwiWSciXSxbNiwyLCJCIl0sWzcsMSwiWiJdLFs4LDAsImV4dChaKSJdLFs0LDMsIl57XFx1bGNvcm5lcn1BK197Zl97ZXh0fTtmLGdfe2V4dH07Z31CIl0sWzAsMV0sWzEsMl0sWzMsMiwiZiJdLFs0LDMsImZfe2V4dH0iXSxbNCw1LCJnX3tleHR9IiwyXSxbNSw2LCJnIiwyXSxbNyw2XSxbOCw3XSxbMiw5XSxbNiw5XV0=
\[
\trimbox{0cm 0cm 0cm 0.75cm}{
\adjustbox{scale=0.9,center}
{\tikzfig{./figures/monoidal_egraphs/pushout_e_cospans}}
}
\]
This pushout exists in $\catname{EHyp}_{\mathrm{s}}(\Sigma)$
by~\cref{prop:pushout_exists}: the apex $m$ is pointed discrete, and both legs
are strict by the third condition on extended cospans, so in particular one of
them is strict and the other anchored.
Then, the result of composition is the extended cospan
\[
n \xrightarrow{f_{\mathrm{ext}} \fatsemi \iota_{1}} n' \sqcup (m'' \setminus m) \xrightarrow{h_{1}} \mathcal{H} \xleftarrow{h_{2}} k' \sqcup (m' \setminus m) \xleftarrow{g'_{\mathrm{ext}} \fatsemi \iota_{1}} k~,
\]
where $\iota_{1}$ is the first coprojection of $\sqcup$ and, writing
$\restriction$ for the restriction of a map to a sub-e-hypergraph of its
pointed discrete domain,
\[
    h_{1} = [\, f_{\mathrm{int}} \fatsemi p_{1}, ~(g_{\mathrm{int}} \fatsemi p_{2}) \restriction_{m'' \setminus m} \,]~,
\qquad
    h_{2} = [\, g'_{\mathrm{int}} \fatsemi p_{2}, ~(f'_{\mathrm{int}} \fatsemi p_{1}) \restriction_{m' \setminus m} \,]~.
\]

The identity on $n$ and the monoidal structure are inherited
from~\cref{def:discrete_functor}.
The functor $D$ extends to
$\widehat{D} : \mathbb{F} \to \catname{EHypI}_{\bullet}(\Sigma)$, sending
$[n]$ to $D[n]$ and $\varphi : [n] \to [m]$ to the extended cospan
\[
D[n] \xrightarrow{\;\id\;} D[n] \xrightarrow{\;D(\varphi)\;} D[m]
      \xleftarrow{\;\id\;} D[m] \xleftarrow{\;\id\;} D[m]~,
\]
whose external legs are identities and which therefore has no strictly internal
interface, so that the conditions on extended cospans and well-typedness hold
vacuously.
The identity on $n$ is $\widehat{D}(\id)$; the monoidal product is $\sqcup$ on
objects and componentwise on extended cospans, with unit $\mathfrak{T} = D[0]$;
and the symmetries are $\sym_{n,m} \defeq \widehat{D}(\sigma_{n,m})$, where
$\sigma_{n,m} : [n+m] \to [m+n]$ is the block transposition.
The symmetric monoidal axioms then follow from those of $(\mathbb{F}, +, [0])$,
since $\widehat{D}$ is strict monoidal by~\cref{remark:discrete_functor_strict}.
\end{definition}

% \begin{remark}[Comparison with~\cite{tiurin2025equivalencehypergraphsdporewriting}]\label{remark:compare_ehypi}
% The extended cospans of~\cite[Defs.~IV.5--IV.6]{tiurin2025equivalencehypergraphsdporewriting}
% have unrooted discrete feet, with two side conditions --- the images of
% external interfaces are top-level, those of strictly internal interfaces are
% not --- and their isomorphisms permute internal interface vertices only
% within the blocks induced by pulling the consistency relation of the carrier
% back along the internal legs.
% Composition there glues carriers by a pushout along the shared external
% foot, and the monoidal product is the disjoint union with the empty
% e-hypergraph as unit.
% The present definition departs from this as follows.
% All four legs are morphisms of the single category
% $\catname{EHyp}_{\leq}(\Sigma)$, with the external composites strict and the
% internal legs lax; the two side conditions become conditions (3)--(4),
% phrased through the root.
% The blocks governing isomorphism are induced by shared immediate parents of
% images, which agrees with the pulled-back consistency relation on strictly
% internal vertices by \cref{remark:compare_ehypergraph}; external vertices,
% previously forming singleton blocks, now jointly form the
% $\top_{\mathcal{G}}$-block, on which order preservation is in any case
% forced by commutativity with the fixed feet, so the isomorphism classes are
% unchanged.
% The pushout defining composition is taken along the \emph{pointed} foot, so
% that the two roots are identified, and the monoidal product is accordingly
% the coslice coproduct $\sqcup$ with unit $T$.
% Finally, the span of the composition pushout lies in
% $\catname{EHyp}_{\mathrm{s}}(\Sigma)$, and lax legs never occur in a pushout
% span: they enter the composite only through $h_{1}$ and $h_{2}$, as
% composites of lax legs with strict injections, which are again lax.
% This division of labour is essential.
% Pushouts along lax legs need not exist: a foot vertex sent to a child of the
% root on one side and into a component of a box on the other would give the
% resulting quotient class two incomparable immediate parents.
% It is precisely this pathology that the pushout construction
% of~\cite[Thm.~A.13]{tiurin2025equivalencehypergraphsdporewriting} guards against with its
% hypotheses --- that interface images share parent sets and consistency
% classes, and that at most one leg assigns parents at all --- and handles in
% its construction by clauses inferring the parent of a parentless element
% from elements it is connected to by a path.
% For spans of anchored legs out of a pointed discrete apex, the hypotheses
% hold automatically --- all images are children of a single anchor, in a
% single component --- and the inference clauses are not needed, as parentless
% elements other than roots do not exist and the roots are identified through
% the foot.
% \end{remark}

As noted in~\cref{example:non_isomorphic_cospans}, the isomorphism of extended cospans does not capture all equalities between morphisms in $\mathbf{S}^{\join}$.
To establish the equivalence between $\mathbf{S}^{\join}$ and $\catname{EHypI}_{\bullet}(\Sigma)$ and moreover establish the equivalence of rewriting in both, we will need to quotient the morphisms (i.e., cospans) of the latter.
For this, we will need to introduce a rewriting framework for extended cospans which we abusively call \emph{DPOI} rewriting.

\section{DPOI rewriting for e-hypergraphs}

Because extended cospans have a more general notion of interface,  including \textit{internal} vertices,  DPOI rewriting as presented in~\cref{def:dpoi_rewriting} needs some adjustments.
We do not expect internal interfaces to be preserved during rewriting: replacing $\boxsort$-edge with another $\boxsort$-edge may arbitrary change the internal interface.
Thus,  we wish for a rewrite rule to be a pair of \textit{extended} cospans of e-hypergraphs with matching \textit{external} (but not necessarily internal) interfaces,  as follows: 
\[
    n \xrightarrow{} n' \xrightarrow{} \mathcal{L} \xleftarrow{} m' \xleftarrow{} m,
\qquad
    n \xrightarrow{} n'' \xrightarrow{} \mathcal{R} \xleftarrow{} m'' \xleftarrow{} m.
\] 
Analogously to~\cref{remark:dpoi_interface_swap}, observing that the following extended cospans express the same data:
\[
    0 \xrightarrow{} 0 \xrightarrow{} \mathcal{L} \xleftarrow{} n'' \sqcup m'' \xleftarrow{} n \sqcup m
\]
\[
    0 \xrightarrow{} 0 \xrightarrow{} \mathcal{R} \xleftarrow{} n' \sqcup m' \xleftarrow{} n \sqcup m
\] allows us to encode rewrite rules as extended cospans of the following form:
\[
\mathcal{L} \xleftarrow{} n' \sqcup m' \xleftarrow{} n \sqcup m \xrightarrow{} n'' \sqcup m'' \xrightarrow{} \mathcal{R},
\] which will fit into the DPO formalism. We make this identification throughout,  without further mention. 

We will also need to restrict the set of candidate sub-e-hypergraphs that will constitute a \emph{match}.

% \begin{definition}[Sub-e-hypergraph]\label{def:sub_ehypergraph}
% Let $\mathcal{G}$ be a rooted e-hypergraph.
% A \emph{sub-e-hypergraph} of $\mathcal{G}$ is a pair of subsets
% $V' \subseteq V_{\mathcal{G}}$, $E' \subseteq E_{\mathcal{G}}$ together with a
% distinguished $r \in E'$ such that
% \begin{enumerate}
%   \item for every $e \in E'$, every vertex occurring in $s_{\mathcal{G}}(e)$ or
%         $t_{\mathcal{G}}(e)$ lies in $V'$;
%   \item $l_{\mathcal{G}}(r) = \cmpsort$, and every
%         $x \in (V' \coprod E') \setminus \{\kappa^{E}(r)\}$ satisfies
%         ${<^{\mu}_{\mathcal{G}}}(x) \in E'$;
%   \item every $\boxsort$-edge in $E'$ has at least two immediate children in
%         $E'$.
% \end{enumerate}
% We regard it as the rooted e-hypergraph
% $(V', E', s_{\mathcal{G}} \restriction_{E'}, t_{\mathcal{G}} \restriction_{E'},
% l_{\mathcal{G}} \restriction_{E'}, r, {<^{\mu}_{\mathcal{G}}} \restriction)$,
% and call $r$ its root.
% \end{definition}

\begin{remark}
Recall the definition of a sub-e-hypergraph from~\cref{def:sub_ehypergraph}.
Equivalently, a sub-e-hypergraph $\mathcal{F}$ of $\mathcal{G}$ is given by a monic homomorphism $\phi : \mathcal{F} \to \mathcal{G}$.
\end{remark}

\begin{lemma}\label{lemma:sub_ehypergraph_image}
The inclusion of a sub-e-hypergraph is an anchored mono with anchor $r$, and it
is strict exactly when $r = \top_{\mathcal{G}}$.
Conversely, if $\phi : \mathcal{F} \to \mathcal{G}$ is an anchored mono, then
$(\phi_{V}(V_{\mathcal{F}}), \phi_{E}(E_{\mathcal{F}}), \phi_{E}(\top_{\mathcal{F}}))$
is a sub-e-hypergraph of $\mathcal{G}$, isomorphic to $\mathcal{F}$.
\end{lemma}

\begin{definition}[Down-closed sub-e-hypergraph]\label{def:down_closed}
A sub-e-hypergraph $\mathcal{F}$ of $\mathcal{G}$ with root $r$ is
\emph{down-closed} if for every $e \in E_{\mathcal{F}}$ with $e \neq r$ and
every $x \in V_{\mathcal{G}} \coprod E_{\mathcal{G}}$ with
$e <^{\mu}_{\mathcal{G}} x$, we have $x \in V_{\mathcal{F}} \coprod E_{\mathcal{F}}$.
\end{definition}

The above definition ensures that a down-closed sub-e-hypergraph is closed under successors of edges other than $r$.

\begin{example}
Consider the sub-e-hypergraph given by a homomorphism $\phi$ below
\[
\adjustbox{max width=0.75\textwidth}{
  \tikzfig{./figures/monoidal_egraphs/non_down_closed_example}
}~.
\]
This sub-e-hypergraph is not down-closed, since the image does not include all the successors of the $\boxsort$-edge $E_3$, in particular the edge $E_2'$.
On the other hand, it is completely fine to not include the edge $h$ and its target with are the successors of the $\top$ in the sub-e-hypergraph---it does not affect the down-closedness of the latter.
Finally note that we are forced to map the edge $f$ together with its source and target in the left e-hypergraph to the edge $f$ in the $E_2$ component edge of the right e-hypergraph: it is forced on us by $\phi$ being a lax homomorphism, mapping it to $E_2'$ component instead would violate the preservation of the parent relation.
\end{example}

% \begin{definition}[Pushout complement]\label{def:pushout_complement}
% Let
% \[
% \mathcal{L} \xleftarrow{} i' \sqcup j' \xleftarrow{} i \sqcup j
% \]
% be an extended cospan, write $f : i \sqcup j \to \mathcal{L}$ for the composite
% of its two legs, and let $m : \mathcal{L} \to \mathcal{G}$ be an anchored mono.
% A \emph{pushout complement} of $(f, m)$ is a pair
% $i \sqcup j \xrightarrow{c} \mathcal{C} \xrightarrow{d} \mathcal{G}$ of
% morphisms of $\catname{EHyp}_{\bullet}(\Sigma)$, with $c$ anchored, such that
% $f \fatsemi m = c \fatsemi d$ and the square
% % https://q.uiver.app/#q=WzAsNCxbMCwwLCJpIFxcc3FjdXAgaiJdLFsyLDAsIlxcbWF0aGNhbHtMfSJdLFswLDIsIlxcbWF0aGNhbHtDfSJdLFsyLDIsIlxcbWF0aGNhbHtHfSJdLFswLDEsImYiXSxbMCwyLCJjIiwyXSxbMSwzLCJtIl0sWzIsMywiZCIsMl0sWzMsMCwiIiwyLHsic3R5bGUiOnsibmFtZSI6ImNvcm5lciJ9fV1d
% \[\begin{tikzcd}
% 	{i \sqcup j} && {\mathcal{L}} \\
% 	\\
% 	{\mathcal{C}} && {\mathcal{G}}
% 	\arrow["f", from=1-1, to=1-3]
% 	\arrow["c"', from=1-1, to=3-1]
% 	\arrow["m", from=1-3, to=3-3]
% 	\arrow["d"', from=3-1, to=3-3]
% 	\arrow["\lrcorner"{anchor=center, pos=0.125, rotate=180}, draw=none, from=3-3, to=1-1]
% \end{tikzcd}\]
% is a pushout in $\catname{EHyp}_{\bullet}(\Sigma)$.
% \end{definition}

% \begin{remark}\label{remark:complement_hypotheses}
% The span $\mathcal{L} \xleftarrow{f} i \sqcup j \xrightarrow{c} \mathcal{C}$
% satisfies the hypotheses of~[the pushout proposition]: its apex is pointed
% discrete, $f$ is strict by the third condition on extended cospans, and $c$ is
% anchored.
% The square above is therefore a pushout precisely when the comparison
% morphism $\mathcal{L} \coprod_{i \sqcup j} \mathcal{C} \to \mathcal{G}$ induced
% by $m$ and $d$ is an isomorphism, so a pushout complement exhibits
% $\mathcal{G}$ as $\mathcal{L}$ glued to $\mathcal{C}$ along the rule's external
% interface.
% Note also that the ambient category here is
% $\catname{EHyp}_{\bullet}(\Sigma)$, in which both $c$ and $d$ are morphisms,
% even though the interface legs of the surrounding extended cospans are merely
% lax.
% \end{remark}

\begin{definition}[Extended pushout complement]\label{def:extended_complement}
Let
\[
\mathcal{L} \xleftarrow{f_{\mathrm{int}}} i' \sqcup j' \xleftarrow{f_{\mathrm{ext}}} i \sqcup j
\qquad\text{and}\qquad
\mathcal{G} \xleftarrow{g_{\mathrm{int}}} n' \sqcup k' \xleftarrow{g_{\mathrm{ext}}} n \sqcup k
\]
be well-typed extended cospans, write $f \defeq f_{\mathrm{ext}} \fatsemi f_{\mathrm{int}}$,
and let $m : \mathcal{L} \to \mathcal{G}$ be an anchored homomorphism with
anchor $a \defeq m_{E}(\top_{\mathcal{L}})$.
An \emph{extended pushout complement} of $(f, m)$ is a pushout complement
$i \sqcup j \xrightarrow{c} \mathcal{L}^{\bot} \xrightarrow{l^{\bot}} \mathcal{G}$
of $(f, m)$, as in the diagram below
% https://q.uiver.app/#q=WzAsNyxbMCwwLCJcXG1hdGhjYWx7TH0iXSxbMSwwLCJpJyBcXGNvcHJvZCBqJyJdLFsyLDAsImkgXFxjb3Byb2QgaiJdLFswLDIsIlxcbWF0aGNhbHtHfSJdLFsyLDIsIlxcbWF0aGNhbHtMfV57XFxib3R9Il0sWzIsNCwibiBcXGNvcHJvZCBrIl0sWzEsMywibicgXFxjb3Byb2QgayciXSxbMCwzLCJtIiwyXSxbMSwwXSxbMiwxXSxbMiw0LCJjIl0sWzQsM10sWzUsNl0sWzYsM10sWzUsNCwiZCIsMl0sWzMsMiwiIiwxLHsic3R5bGUiOnsibmFtZSI6ImNvcm5lciJ9fV1d
\[\begin{tikzcd}
	{\mathcal{L}} & {i' \coprod j'} & {i \coprod j} \\
	\\
	{\mathcal{G}} && {\mathcal{L}^{\bot}} \\
	& {n' \coprod k'} \\
	&& {n \coprod k}
	\arrow["m"', from=1-1, to=3-1]
	\arrow[from=1-2, to=1-1]
	\arrow[from=1-3, to=1-2]
	\arrow["c", from=1-3, to=3-3]
	\arrow["\lrcorner"{anchor=center, pos=0.125, rotate=90}, draw=none, from=3-1, to=1-3]
	\arrow[from=3-3, to=3-1]
	\arrow[from=4-2, to=3-1]
	\arrow["d"', from=5-3, to=3-3]
	\arrow[from=5-3, to=4-2]
\end{tikzcd}~,\]
such that
\begin{enumerate}
  \item for every $e \in E_{\mathcal{L}}$ with $e \neq \top_{\mathcal{L}}$ and
        every $x \in V_{\mathcal{G}} \coprod E_{\mathcal{G}}$ with
        $m_{E}(e) <^{\mu}_{\mathcal{G}} x$, the element $x$ lies in the image of
        $m$;
  \item writing
        $S \defeq g_{\mathrm{int}}^{-1}\bigl(m(f_{\mathrm{int}}(i' \setminus i))
          \sqcup m(f_{\mathrm{int}}(j' \setminus j))\bigr)$,
        the restriction of $g_{\mathrm{int}}$ to $(n' \sqcup k') \setminus S$
        factors through $l^{\bot}$.
\end{enumerate}
\end{definition}

\begin{remark}\label{remark:complement_cases}
Condition (1) forbids the following matches
\[
\adjustbox{max width=0.75\textwidth}{
  \tikzfig{./figures/monoidal_egraphs/extended_pushout_complement_condition_1}
}~.
\]
Intuitively, the image must contain all successors of every $\boxsort$-edge and hence all $\cmpsort$-edges and their successors if they have a $\boxsort$-edge as a predecessor.
Condition (2) disables the following matches
\[
\adjustbox{max width=0.65\textwidth}{
  \tikzfig{./figures/monoidal_egraphs/extended_pushout_complement_condition_2}
}
\] 
The pictured pushout complement is a usual pushout complement---computing the pushout closes the square, however, it is not an extended pushout complement as it violates the condition (2).
Computing $S = g_{\mathrm{int}}^{-1}\bigl(m(f_{\mathrm{int}}(i' \setminus i)) \coprod m(f_{\mathrm{int}}(j' \setminus j))\bigr)$ yields $v_1, v_2$ portion of the internal interface and the restriction $g_{\mathrm{int}} \restriction S$ does not factor through $l^{\bot}$ as there is no vertex to match with $v_1$.
Put simply, this condition tells us to remove the part of $n' \coprod m'$ that matches the image of the strictly intrnal interface of $\mathcal{L}$---since $\mathcal{L}$ is removed in the pushout complement, its strictly internal interfaces must be removed as well.

The reason for removing only the portion of strictly internal interface can be grasped from the following example.
\[
\adjustbox{max width=0.75\textwidth}{
\tikzfig{./figures/monoidal_egraphs/extended_pushout_complement_condition_3}
}
\]
This time the internal interface of $\mathcal{L}$ is not empty.
Computing $S$ gives us $v_3, v_4, v_5$ and this factors through the depicted pushout complement.
Had we removed the entire internal interface, we would have been left with $v_4, v_5$ with no way to recover that interface if we were to glue something to $v_3$.
\end{remark}

\begin{definition}[EDPOI rewriting]\label{def:dpoi_rooted}
A \emph{rewrite rule} is a span of extended cospans with matching external
interfaces,
\begin{align*}
\mathcal{L} \xleftarrow{} i' \sqcup j' \xleftarrow{} i \sqcup j
  \xrightarrow{} i'' \sqcup j'' \xrightarrow{} \mathcal{R}~. \tag{$*$}
\end{align*}
Given an extended span
\[
\mathcal{G} \xleftarrow{} n' \sqcup k' \xleftarrow{} n \sqcup k
  \xrightarrow{} n'' \sqcup k'' \xrightarrow{} \mathcal{H}~,
\]
we say that $\mathcal{G}$ \emph{rewrites to} $\mathcal{H}$ with external
interface $n \sqcup k$, taking internal interface $n' \sqcup k'$ to
$n'' \sqcup k''$---written $\mathcal{G} \;\overline{\rightsquigarrow}\; \mathcal{H}$---via the rule $(*)$ if
there exist an anchored homomorphism $m : \mathcal{L} \to \mathcal{G}$, an
object $\mathcal{L}^{\bot}$, and morphisms completing the following commutative
diagram
% https://q.uiver.app/#q=WzAsMTEsWzAsMCwiXFxtYXRoY2Fse0x9Il0sWzEsMCwiaScgXFxjb3Byb2QgaiciXSxbMiwwLCJpIFxcY29wcm9kIGoiXSxbMCwyLCJcXG1hdGhjYWx7R30iXSxbMiwyLCJcXG1hdGhjYWx7TH1ee1xcYm90fSJdLFsyLDQsIm4gXFxjb3Byb2QgayJdLFsxLDMsIm4nIFxcY29wcm9kIGsnIl0sWzQsMiwiXFxtYXRoY2Fse0h9Il0sWzQsMCwiXFxtYXRoY2Fse1J9Il0sWzMsMywibicnIFxcY29wcm9kIGsnJyJdLFszLDAsImknJ1xcY29wcm9kIGonJyJdLFswLDMsIm0iLDJdLFsxLDAsImZfe1xcbWF0aHJte2ludH19IiwyXSxbMiwxLCJmX3tcXG1hdGhybXtleHR9fSIsMl0sWzIsNCwiYyJdLFs0LDMsImxee1xcYm90fSIsMl0sWzUsNiwiZ197XFxtYXRocm17ZXh0fX0iXSxbNiwzLCJnX3tcXG1hdGhybXtpbnR9fSJdLFs1LDQsImQiLDJdLFszLDIsIiIsMSx7InN0eWxlIjp7Im5hbWUiOiJjb3JuZXIifX1dLFs0LDcsIlxcd2lkZXRpbGRle1xcaW90YV97XFxtYXRoY2Fse0x9XntcXGJvdH19fSJdLFs4LDcsIlxcd2lkZXRpbGRle1xcaW90YV97XFxtYXRoY2Fse1J9fX0iXSxbNSw5LCJoX3tcXG1hdGhybXtleHR9fSIsMl0sWzksNywiaF97XFxtYXRocm17aW50fX0iLDJdLFs3LDIsIiIsMSx7InN0eWxlIjp7Im5hbWUiOiJjb3JuZXIifX1dLFsyLDEwLCJwX3tcXG1hdGhybXtleHR9fSJdLFsxMCw4LCJwX3tcXG1hdGhybXtpbnR9fSJdXQ==
\[\begin{tikzcd}
	{\mathcal{L}} & {i' \coprod j'} & {i \coprod j} & {i''\coprod j''} & {\mathcal{R}} \\
	\\
	{\mathcal{G}} && {\mathcal{L}^{\bot}} && {\mathcal{H}} \\
	& {n' \coprod k'} && {n'' \coprod k''} \\
	&& {n \coprod k}
	\arrow["m"', from=1-1, to=3-1]
	\arrow["{f_{\mathrm{int}}}"', from=1-2, to=1-1]
	\arrow["{f_{\mathrm{ext}}}"', from=1-3, to=1-2]
	\arrow["{p_{\mathrm{ext}}}", from=1-3, to=1-4]
	\arrow["c", from=1-3, to=3-3]
	\arrow["{p_{\mathrm{int}}}", from=1-4, to=1-5]
	\arrow["{\widetilde{\iota_{\mathcal{R}}}}", from=1-5, to=3-5]
	\arrow["\lrcorner"{anchor=center, pos=0.125, rotate=90}, draw=none, from=3-1, to=1-3]
	\arrow["{l^{\bot}}"', from=3-3, to=3-1]
	\arrow["{\widetilde{\iota_{\mathcal{L}^{\bot}}}}", from=3-3, to=3-5]
	\arrow["\lrcorner"{anchor=center, pos=0.125, rotate=180}, draw=none, from=3-5, to=1-3]
	\arrow["{g_{\mathrm{int}}}", from=4-2, to=3-1]
	\arrow["{h_{\mathrm{int}}}"', from=4-4, to=3-5]
	\arrow["d"', from=5-3, to=3-3]
	\arrow["{g_{\mathrm{ext}}}", from=5-3, to=4-2]
	\arrow["{h_{\mathrm{ext}}}"', from=5-3, to=4-4]
\end{tikzcd}\]
such that the two marked squares are pushouts and the following conditions
hold.
\begin{enumerate}
  \item $i \sqcup j \xrightarrow{c} \mathcal{L}^{\bot} \xrightarrow{l^{\bot}} \mathcal{G}$
        is an extended pushout complement of $(f, m)$, with chosen
        factorisation $g_{1} \sqcup g_{2}$, where
        \[
        S_{n} \defeq g_{\mathrm{int}}^{-1}\bigl(m(f_{\mathrm{int}}(i' \setminus i))\bigr)~,
        \qquad
        S_{k} \defeq g_{\mathrm{int}}^{-1}\bigl(m(f_{\mathrm{int}}(j' \setminus j))\bigr)~,
        \]
        so that $g_{1} : n' \setminus S_{n} \to \mathcal{L}^{\bot}$ and
        $g_{2} : k' \setminus S_{k} \to \mathcal{L}^{\bot}$.
  \item The internal interfaces of $\mathcal{H}$ are
        \[
        n'' = (n' \setminus S_{n}) \sqcup (i'' \setminus i)~,
        \qquad
        k'' = (k' \setminus S_{k}) \sqcup (j'' \setminus j)~.
        \]
  \item The internal legs of $\mathcal{H}$ are
        \[
        h_{\mathrm{int}} \restriction_{n''}
          = [\, g_{1} \fatsemi \widetilde{\iota_{\mathcal{L}^{\bot}}},\;
                (p_{\mathrm{int}} \fatsemi \widetilde{\iota_{\mathcal{R}}}) \restriction_{i'' \setminus i} \,]~,
        \]
        and dually on $k''$, and its external legs are
        $h_{\mathrm{ext}} = g_{\mathrm{ext}}$ followed by the evident inclusion
        of $n \sqcup k$ into $n'' \sqcup k''$.
\end{enumerate}
Given a set $\mathfrak{R}$ of rewrite rules and extended cospans $\mathcal{G}$
and $\mathcal{H}$, we write $\mathcal{G} \;\overline{\rightsquigarrow}_{\mathfrak{R}}\; \mathcal{H}$
if $\mathcal{G}$ rewrites to $\mathcal{H}$ via some rule of $\mathfrak{R}$, and
$\mathcal{G} \;\overline{\rightsquigarrow}^{*}_{\mathfrak{R}}\; \mathcal{H}$ for the reflexive,
symmetric and transitive closure of $\rightsquigarrow_{\mathfrak{R}}$.
\end{definition}

\begin{lemma}\label{lemma:rewrite_yields_ecospan}
Suppose $\mathcal{G} \;\overline{\rightsquigarrow}\; \mathcal{H}$ as in~\cref{def:dpoi_rooted},
via a rule with well-typed extended cospans $\mathcal{L}$ and $\mathcal{R}$ and
a well-typed $\mathcal{G}$.
Then
\[
n \xrightarrow{} n'' \xrightarrow{f_{1}} \mathcal{H}
  \xleftarrow{f_{2}} k'' \xleftarrow{} k
\]
is a well-typed extended cospan.
\end{lemma}
\begin{proof}
  TODO
\end{proof}

\section{Soundness and completeness.}

First, let us define $\join$ of two extended cospans, i.e., morphisms in $\catname{EHypI}_{\bullet}(\Sigma)$.
\begin{definition}[Join of extended cospans]\label{def:join_of_extended_cospans}
Given two extended cospans 
\[
n \xrightarrow{} n' \xrightarrow{} \mathcal{F} \xleftarrow{} m' \xleftarrow{} m
\]
and
\[
n \xrightarrow{} n'' \xrightarrow{} \mathcal{G} \xleftarrow{} m'' \xleftarrow{} m
\]
we define their \emph{join} as the extended cospan of~\cref{fig:join_of_extended_cospans},
whose carrier is obtained from the disjoint union of the underlying vertex and
edge sets of $\mathcal{F}$ and $\mathcal{G}$ by adjoining a fresh root $\top$, a
$\boxsort$-edge $b$ with $\top <^{\mu} b$ and $b <^{\mu} \top_{\mathcal{F}}$,
$b <^{\mu} \top_{\mathcal{G}}$, and fresh vertices $s(b)$ and $t(b)$ with
$|s(b)| = |V_{n}|$ and $|t(b)| = |V_{m}|$, also children of $\top$;
the labels, sources and targets are inherited, with $\top_{\mathcal{F}}$ and
$\top_{\mathcal{G}}$ retaining their label $\cmpsort$ and becoming the two
components of $b$.

Its external input leg sends $n$ isomorphically onto $s(b)$, and dually for $m$
and $t(b)$; its internal input interface is $s(b) \sqcup n' \sqcup n''$, with
the block of $\top_{\mathcal{F}}$ ordered so that its $k$-th vertex is
$f_{\mathrm{ext}}$ of the $k$-th vertex of $n$, and likewise for
$\top_{\mathcal{G}}$ via $g_{\mathrm{ext}}$; dually for the outputs.
This generalises to $k$ summands in the obvious way.
\end{definition}

\begin{figure}
\[
\adjustbox{max width=0.65\textwidth}{
\tikzfig{./figures/monoidal_egraphs/f_plus_g_new}
}
\]
\caption{The join of two extended cospans.}
\label{fig:join_of_extended_cospans}
\end{figure}

\begin{definition}[Structural rules]\label{def:structural_rules}
Let $\mathcal{R}$ be the set of EDPOI rewrite rules $\langle \mathcal{L}, \mathcal{R} \rangle$
and $\langle \mathcal{R}, \mathcal{L} \rangle$ obtained from the equations
of~\cref{fig:string-equations}.
\end{definition}

\begin{definition}[Quotient by rewrites]\label{def:quotient_by_rewrites}
$\catname{EHypI}_{\bullet}(\Sigma) / \mathcal{R}$ is the category with the same
objects as $\catname{EHypI}_{\bullet}(\Sigma)$ and, as morphisms $n \to m$, the
equivalence classes of morphisms under $\overline{\rightsquigarrow}^{*}_{\mathcal{R}}$.
\end{definition}

Let $\mathbf{H}^{\join}_{\Sigma}$ be the Frobenius PROP obtained by applying $\widetilde{F}$ to $\mathbf{H}_{\Sigma}$.
\begin{definition}\label{not:component_cospan}
Let $\mathcal{G}$ be a well-typed extended cospan and $c$ a $\cmpsort$-edge of
its carrier.
We write $\mathcal{G}_{c}$ for the extended cospan whose carrier is the
sub-e-hypergraph consisting of $c$ and all its descendants, with root $c$, whose
external interfaces are the blocks $\pi_{f}^{-1}(c)$ and $\pi_{g}^{-1}(c)$
of~\cref{def:isomorphic_ecospans}, ordered so that their $k$-th elements
correspond to the $k$-th source and target of the $\boxsort$-edge above $c$, and
whose internal interfaces are the restrictions of $n'$ and $m'$ to the vertices
lying strictly below $c$.
In particular $\mathcal{G}_{\top_{\mathcal{G}}} = \mathcal{G}$, and $\Phi(c_{i})$
in~\cref{def:comparison_functor} abbreviates $\Phi(\mathcal{G}_{c_{i}})$.
If $c$ is a component of a $\boxsort$-edge $b$, then $\mathcal{G}_{c}$ has type
$|s(b)| \to |t(b)|$ by well-typedness.
\end{definition}

\begin{lemma}\label{lemma:component_cospan_wf}
$\mathcal{G}_{c}$ is a well-typed extended cospan.
\end{lemma}
\begin{proof}
The carrier satisfies the conditions of~\cref{def:sub_ehypergraph}: condition
(1) holds since a vertex incident to an edge shares that edge's immediate
parent, so it is a descendant of $c$ whenever the edge is; condition (2) holds
with $r = c$ because the set is closed downwards, so every element other than
$c$ retains its immediate parent; and condition (3) is inherited from
non-degeneracy in $\mathcal{G}$, since a $\boxsort$-edge below $c$ keeps all its
children.
The external legs are mono and order-preserving by construction.
Their composites are strict: by~\cref{def:isomorphic_ecospans} the block
$\pi_{f}^{-1}(c)$ consists precisely of the interface vertices whose immediate
parent is $c$, which is the root of $\mathcal{G}_{c}$; the same fact shows that
the vertices of the internal interface lying strictly below $c$ are not children
of $c$, so the strictly internal interface of $\mathcal{G}_{c}$ is disjoint from
its external one.
Well-typedness of each $\boxsort$-edge below $c$ is the corresponding instance
of well-typedness of $\mathcal{G}$, the blocks being unchanged.
\end{proof}

\begin{definition}\label{def:comparison_functor}
Define $\Phi : \catname{EHypI}_{\bullet}(\Sigma) \to \mathbf{H}^{\join}_{\Sigma}$
by $\Phi(n) = |V_{n}|$ on objects and, on a well-typed extended cospan
$\mathcal{G}$, by recursion on the depth of its hierarchy: $\Phi(\mathcal{G})$
is the image of $\mathcal{G}^{\flat}$ under the isomorphism
$\catname{HypI}(\Sigma) \cong \mathbf{H}_{\Sigma}$, with each box generator $b$
interpreted as $\sum_{i} \Phi(c_{i})$, the join over the components of $b$.
\end{definition}

\begin{lemma}\label{lemma:phi_well_defined}
$\Phi$ is constant on isomorphism classes of extended cospans.
\end{lemma}
\begin{proof}
By induction on the depth of the hierarchy,
using~\cref{def:depth}.
Let $(\alpha, \beta, \gamma)$ be an isomorphism as
in~\cref{def:isomorphic_ecospans}.
Since $\beta$ is an isomorphism of $\catname{EHyp}_{\mathrm{s}}(\Sigma)$ it is
strict, so it sends $\top_{\mathcal{G}}$ to $\top_{\mathcal{G}'}$, preserves
$<^{\mu}$, and preserves labels; hence it restricts to a bijection between the
children of the two roots and between their $\boxsort$-edges, inducing an
isomorphism $\mathcal{G}^{\flat} \cong \mathcal{G}'^{\flat}$ under which a
$\boxsort$-generator $b$ corresponds to $\beta(b)$.

Next, $\alpha$ maps blocks to blocks: for $x \in n'$,
\[
\pi_{h}(\alpha(x))
= {<^{\mu}_{\mathcal{G}'}}\bigl(h_{\mathrm{int}}(\alpha(x))\bigr)
= {<^{\mu}_{\mathcal{G}'}}\bigl(\beta(f_{\mathrm{int}}(x))\bigr)
= \beta\bigl({<^{\mu}_{\mathcal{G}}}(f_{\mathrm{int}}(x))\bigr)
= \beta(\pi_{f}(x))~,
\]
using commutativity of the diagram and strictness of $\beta$; so $\alpha$
restricts to a bijection from the block of $c$ to the block of $\beta(c)$, and
by~\cref{def:isomorphic_ecospans} this restriction is order-preserving whenever
$c \neq \top_{\mathcal{G}}$.
Therefore $(\alpha, \beta, \gamma)$ restricts to an isomorphism
$\mathcal{G}_{c} \cong \mathcal{G}'_{\beta(c)}$ for every $\cmpsort$-edge
$c \neq \top_{\mathcal{G}}$, and by the inductive hypothesis
$\Phi(\mathcal{G}_{c}) = \Phi(\mathcal{G}'_{\beta(c)})$.
The two flattened cospans thus have equal images under the interpretation of
their $\boxsort$-generators, so $\Phi(\mathcal{G}) = \Phi(\mathcal{G}')$.
\end{proof}

\begin{lemma}\label{lemma:phi_functor}
$\Phi$ is a strict symmetric monoidal functor.
\end{lemma}
\begin{proof}
Write $B_{\mathcal{G}}$ for the set of $\boxsort$-edges of $\mathcal{G}$ that are
children of $\top_{\mathcal{G}}$ and $\Sigma_{\mathcal{G}}$ for $\Sigma$ extended
by a generator $\sigma_{b} : |s(b)| \to |t(b)|$ for each $b \in B_{\mathcal{G}}$,
so that $\mathcal{G}^{\flat}$ is a morphism of
$\catname{HypI}(\Sigma_{\mathcal{G}})$ and $\Phi(\mathcal{G})$ is its image under
the PROP morphism $\mathbf{H}_{\Sigma_{\mathcal{G}}} \to \mathbf{H}^{\join}_{\Sigma}$
which fixes $\Sigma$ and sends $\sigma_{b}$ to
$\sum_{b <^{\mu} c} \Phi(\mathcal{G}_{c})$.

Composition in~\cref{def:rooted_ehypi} is computed by a pushout along a pointed
discrete foot.
By the construction in the proof of~\cref{prop:pushout_exists} the only edges
identified are the two roots, and no immediate parent below the root is changed;
hence $B_{\mathcal{F} \fatsemi \mathcal{G}} = B_{\mathcal{F}} \sqcup B_{\mathcal{G}}$,
each $\boxsort$-edge keeps its components, $(\mathcal{F} \fatsemi \mathcal{G})_{c}
= \mathcal{F}_{c}$ or $\mathcal{G}_{c}$ accordingly, and
$(\mathcal{F} \fatsemi \mathcal{G})^{\flat} = \mathcal{F}^{\flat} \fatsemi \mathcal{G}^{\flat}$
in $\catname{HypI}(\Sigma_{\mathcal{F}} \cup \Sigma_{\mathcal{G}})$.
Since $\catname{HypI}(-) \cong \mathbf{H}_{(-)}$ is functorial and the PROP
morphism above restricts on each summand to the one used for $\mathcal{F}$ and
for $\mathcal{G}$, we get
$\Phi(\mathcal{F} \fatsemi \mathcal{G}) = \Phi(\mathcal{F}) \fatsemi \Phi(\mathcal{G})$.
The argument for $\otimes$ is the same.
Identities and symmetries are $\widehat{D}$-images
by~\cref{def:discrete_functor,remark:discrete_functor_strict}, hence have no
$\boxsort$-edges, so $\Phi$ sends them to identities and symmetries; and $\Phi$
is strict on objects because $|V_{n \sqcup m}| = |V_{n}| + |V_{m}|$.
\end{proof}

\begin{lemma}\label{lemma:phi_join}
$\Phi(\mathcal{F} + \mathcal{G}) = \Phi(\mathcal{F}) + \Phi(\mathcal{G})$, and
similarly for $k$-ary joins.
\end{lemma}
\begin{proof}
By~\cref{def:join_of_extended_cospans} the children of the root of
$\mathcal{F} + \mathcal{G}$ are exactly the $\boxsort$-edge $b$ and the vertices
$s(b) \sqcup t(b)$, so $(\mathcal{F} + \mathcal{G})^{\flat}$ is $\sigma_{b}$ with
identity wiring and
\[
\Phi(\mathcal{F} + \mathcal{G})
= \Phi\bigl((\mathcal{F}+\mathcal{G})_{\top_{\mathcal{F}}}\bigr)
+ \Phi\bigl((\mathcal{F}+\mathcal{G})_{\top_{\mathcal{G}}}\bigr)~.
\]
The ordering of the blocks fixed in~\cref{def:join_of_extended_cospans} makes
the block of $\top_{\mathcal{F}}$ correspond to $s(b)$ in the order induced by
$\mathcal{F}$'s external leg, so
$(\mathcal{F}+\mathcal{G})_{\top_{\mathcal{F}}} = \mathcal{F}$
in the sense of~\cref{not:component_cospan}, and likewise for $\mathcal{G}$.
\end{proof}

\begin{definition}\label{def:expansion_weight}
For a rooted e-hypergraph $\mathcal{G}$ define $\nu$ on component and box edges
by mutual recursion:
\[
\nu(c) = \bigl(a(c) + 2\bigr) \prod_{c \,<^{\mu}\, b,\; b \in E_{\boxsortsubscript}} \nu(b)~,
\qquad
\nu(b) = 1 + \sum_{b \,<^{\mu}\, c,\; c \in E_{\cmpsortsubscript}} \nu(c)~,
\]
where $a(c)$ is the number of vertices and labelled edges that are immediate
children of $c$, and the empty product is $1$.
Set $\nu(\mathcal{G}) \defeq \nu(\top_{\mathcal{G}})$.
\end{definition}

\begin{lemma}\label{lemma:normal_form}
Every morphism of $\catname{EHypI}_{\bullet}(\Sigma)$ rewrites under
$\mathcal{R}$, applied left to right, to one whose carrier is either box-free or
has a single box edge, a child of the root, whose components are box-free and
pairwise non-isomorphic.
\end{lemma}

\begin{theorem}\label{thm:frobenius_equivalence}
$\mathbf{H}^{\join}_{\Sigma} \simeq \catname{EHypI}_{\bullet}(\Sigma) / \mathcal{R}$.
\end{theorem}


\begin{lemma}[restatement of~\cref{lemma:normal_form}]\label{lemma:normal_form_precise}
Every morphism of $\catname{EHypI}_{\bullet}(\Sigma)$ rewrites under
$\mathcal{R}$, applied left to right, to one whose carrier either has no
$\boxsort$-edge, or has exactly one $\boxsort$-edge $b$, a child of the root,
such that the children of the root are exactly $b$ and the vertices
$s(b) \sqcup t(b)$, and the $\mathcal{G}_{c}$ for the components $c$ of $b$ have
no $\boxsort$-edges and are pairwise non-isomorphic.
\end{lemma}
\begin{proof}
Every non-trivial left-to-right application of a rule
of~\cref{def:structural_rules} strictly decreases $\nu$
of~\cref{def:expansion_weight}, a positive integer.

\emph{Distribution below the root.}
Let $b$ have components $c_{1}, \ldots, c_{p}$ and $c <^{\mu} b$ with
$c \neq \top$, and put $X \defeq (a(c)+2)\prod_{b' \neq b} \nu(b')$, so that
$\nu(c) = X(1 + \sum_{i} \nu(c_{i}))$.
The rule replaces $c$, among the children of its $\boxsort$-edge, by components
$c^{(i)}$ carrying the remaining children of $c$ together with the children of
$c_{i}$, whence
$\nu(c^{(i)}) = \tfrac{a(c)+a(c_{i})+2}{a(c_{i})+2} \cdot \tfrac{X}{a(c)+2} \cdot \nu(c_{i})$.
Since $2(a(c)+a(c_{i})+2) \leq (a(c)+2)(a(c_{i})+2)$, we get
$\sum_{i}\nu(c^{(i)}) \leq X \sum_{i} \nu(c_{i}) < \nu(c)$.

\emph{Distribution at the root.}
With $c = \top$, the rule yields a root whose only child is a $\boxsort$-edge
$B$ with components $c^{(i)}$, so $\nu$ becomes
$2(1 + \sum_{i}\nu(c^{(i)})) \leq 2 + X\sum_{i}\nu(c_{i})$, against
$X + X\sum_{i}\nu(c_{i})$ before; this is strict unless $X = 2$, that is unless
$b$ and its ports are the only children of the root, in which case the rule does
not apply.

\emph{Associativity and idempotence.}
If $b$ is the only child of a component $c$ with $a(c) = 0$, merging $b$ into
$c$'s $\boxsort$-edge replaces the summand $\nu(c) = 2 + 2\sum_{i}\nu(c_{i})$ by
$\sum_{i}\nu(c_{i})$; deleting a duplicate component removes a positive summand.

As $\nu$ is built from sums and products with positive coefficients, a strict
decrease in a subterm is a strict decrease of $\nu(\mathcal{G})$.
The process terminates, and a terminal cospan admits no distribution, no merge
and no duplicate component, which is the stated form.
\end{proof}

\begin{proof}[Proof of~\cref{thm:frobenius_equivalence}]
By~\cref{lemma:phi_well_defined,lemma:phi_functor} $\Phi$ is a functor, and
by~\cref{lemma:phi_sound} it identifies
$\overline{\rightsquigarrow}_{\mathcal{R}}$-related morphisms, so it descends to
$\overline{\Phi} : \catname{EHypI}_{\bullet}(\Sigma)/\mathcal{R} \to
\mathbf{H}^{\join}_{\Sigma}$ by~\cref{def:quotient_by_rewrites}.

\emph{Essential surjectivity.}
$\Phi(D[k]) = k$ by~\cref{def:discrete_functor}, and every object with $k$
vertices is isomorphic to $D[k]$.

\emph{Fullness.}
By~\cref{prop:prop-adjunction} the free semilattice on a set is its set of
non-empty finite subsets, so a morphism of $\mathbf{H}^{\join}_{\Sigma}(n,m)$ is
a non-empty finite set $\{h_{1}, \ldots, h_{p}\}$ of morphisms of
$\mathbf{H}_{\Sigma}(n,m)$.
For $p = 1$ take the cospan corresponding to $h_{1}$ under
$\catname{HypI}(\Sigma) \cong \mathbf{H}_{\Sigma}$; it has no $\boxsort$-edges,
so its image is $\{h_{1}\}$.
For $p \geq 2$ take the join of those $p$ cospans
as in~\cref{def:join_of_extended_cospans}: it is well typed since the $h_{i}$
are parallel, non-degenerate since $p \geq 2$, and has image
$\{h_{1}, \ldots, h_{p}\}$ by~\cref{lemma:phi_join}.

\emph{Faithfulness.}
Suppose $\Phi(\mathcal{G}) = \Phi(\mathcal{H})$.
By~\cref{lemma:normal_form_precise} they rewrite to normal forms $\mathcal{N}$,
$\mathcal{N}'$, which have equal images by~\cref{lemma:phi_sound}.
On cospans without $\boxsort$-edges $\Phi$ is the bijection
$\catname{HypI}(\Sigma) \cong \mathbf{H}_{\Sigma}$ followed by the singleton
map, so such a normal form has a singleton image, whereas one with a
$\boxsort$-edge has image of cardinality its number of components, at least two,
those components being pairwise non-isomorphic and hence, by the same bijection,
of pairwise distinct image.
So either both are $\boxsort$-edge-free and $\mathcal{N} \cong \mathcal{N}'$ by
that bijection; or both have one $\boxsort$-edge, with components
$\mathcal{N}_{1}, \ldots, \mathcal{N}_{p}$ and
$\mathcal{N}'_{1}, \ldots, \mathcal{N}'_{q}$, and by~\cref{lemma:phi_join}
together with the root having no other children,
$\{\Phi(\mathcal{N}_{i})\}_{i} = \{\Phi(\mathcal{N}'_{j})\}_{j}$ with both
families pairwise distinct.
Hence $p = q$ and there is a permutation $\sigma$ with
$\Phi(\mathcal{N}_{i}) = \Phi(\mathcal{N}'_{\sigma(i)})$, so
$\mathcal{N}_{i} \cong \mathcal{N}'_{\sigma(i)}$.
Since the components of a $\boxsort$-edge carry no order,
by~\cref{def:isomorphic_ecospans} $\mathcal{N} \cong \mathcal{N}'$, and
therefore $\mathcal{G} \;\overline{\rightsquigarrow}^{*}_{\mathcal{R}}\; \mathcal{H}$.
\end{proof}

\paragraph{Related work.}
The rooted e-hypergraphs of~\cref{def:rooted_ehypergraph} belong to a broad
family of hierarchical graph formalisms, but their hierarchy has a more
specific purpose than merely decomposing a large graph.
A box edge denotes a join, its component-edge children denote the summands of
that join, and every ordinary edge and vertex is owned by exactly one
component.
Thus, alternation and incidence in \cref{def:rooted_ehypergraph} are
representation invariants for semilattice-enriched string diagrams rather than
optional well-formedness conditions.
This distinction prevents us from adopting the following existing formalisms
unchanged.

Drewes, Hoffmann and Plump define a hierarchical graph recursively as a
top-level graph in which distinguished frame edges each contain one further
hierarchical graph~\cite[Defs.~2.1--2.2]{plump:hierarchical-graphs}.
This strict containment is well suited to abstraction: the contents of
different frames are disjoint, morphisms act level by level and preserve
hierarchical depth, and rule schemata with variables can copy or delete a frame
together with its contents~\cite[Sec.~4]{plump:hierarchical-graphs}.
It does not, however, distinguish several contents of one frame as alternative
representatives of the same morphism.
Putting every alternative in its own frame would not record that the frames
together form one semilattice join, whereas putting the alternatives together
in the contents of one frame would forget where one summand ends and another
begins.
Recovering this information would require component markers, a common scoped
interface, and morphisms for internal interface legs that need not preserve
depth---precisely the additional structure built into rooted e-hypergraphs and
their extended cospans.

Palacz's formalism is the closest at the level of carriers.
It adds to an arbitrary flat graph an acyclic parent-assignment function on all
vertices and edges, together with a formal root, and its morphisms preserve the
parent assignment~\cite[Defs.~1 and~4]{palacz:hierarchical-transform}.
Its generality is also why it cannot be used directly here: any graph atom may
be the parent of any other, and incidence and hierarchy are independent, so an
edge may cross component boundaries.
Consequently, the parent tree by itself neither identifies which levels are
joins and which are their summands nor determines the ownership of wires.
Root-level and nesting morphisms resemble our strict and anchored
homomorphisms, respectively, but Palacz has no counterpart of the lax,
ancestor-preserving maps required for internal interface legs.
Rooted e-hypergraphs may therefore be regarded as a typed, incidence-local
specialisation of Palacz's carriers, but obtaining the required category means
restricting the objects by alternation, component structure and non-degeneracy,
and replacing its single class of morphisms by the three classes in
\cref{def:rooted_homomorphism}.

Castelnovo, Gadducci and Miculan take a different, explicitly categorical
route.
Their hierarchical hypergraphs with interface consist of a tree order on the
set of hyperedges, an otherwise independent incidence map, and a map from an
interface set to the vertices; the resulting category is
adhesive~\cite[Def.~3.5 and Thm.~3.4]{castelnovoNewCriterion}.
This gives a powerful generic setting for DPO rewriting, but only hyperedges
participate in the hierarchy.
It therefore does not record which component owns a vertex, nor does its tree
order distinguish boxes from alternative components.
Moreover, adhesivity of the ambient category does not automatically descend to
the subcategory cut out by our alternation, incidence and non-degeneracy
conditions.
The specialised pushouts needed for extended-cospan composition and DPOI
rewriting must consequently be established for those invariants.

In summary, the earlier constructions supply, respectively, recursive frames,
a general parent-assignment representation, and an adhesive ambient category.
None simultaneously represents a family of alternatives with a shared scoped
interface and provides the strict, anchored and lax maps needed for composition
and rewriting.
The additional structure of rooted e-hypergraphs is what makes the hierarchy a
concrete syntax for joins in a semilattice-enriched symmetric monoidal category,
rather than a generic graph-decomposition mechanism.